\begin{textsample}{POS Dim 3 – 2019\_c6 – Score 16.00 – 2019\_c6/enade\_2019\_tec\_est\_cos\_b618aa5e88.txt}  \label{ex:f3_pos_011}
CAPÍTULO 6 ANÁLISE TÉCNICA DA PROVA Este capítulo tem por objetivo apresentar o desempenho dos estudantes concluintes de Tecnologia em Estética e Cosmética no Enade/2019.
Para isso, foram calculadas as estatísticas básicas da prova em seu todo ( seção 6.1.1 ), bem como as estatísticas relacionadas aos Componentes de Formação Geral ( seção 6.1.2 ) e de Conhecimento Específico da Área ( seção 6.1.3 ).
Dadas as suas características, foram analisadas, em separado, as questões objetivas ( seção 6.2 ) e as questões discursivas ( seção 6.3 ).
Tomando -se como base as duas questões discursivas do Componente de Formação Geral, nas seções 6.3.1.5 a 6.3.1.8, são apresentados comentários sobre a correção das respostas em relação à Língua Portuguesa.
Para as questões objetivas, foram disponibilizados os Índices de Facilidade e de discriminação Ponto-Bisserial, também em separado, para os Componentes de Formação Geral ( seção 6.2.1 ) e de Conhecimento Específico ( seção 6.2.2 ).
De cada componente, uma das questões foi escolhida para exemplificar a análise gráfica, relacionando -se as alternativas escolhidas pelos estudantes ( inclusive o gabarito ) com o número de acertos no componente.
No Anexo I, apresenta -se a íntegra da análise gráfica para todas as questões objetivas.
Para cada uma das questões discursivas, os conteúdos dos tipos mais comuns de respostas dos estudantes são apresentados e comparados com o padrão de resposta esperado ( ver Anexo VIII com o padrão de respostas ).
Nas tabelas, constam as seguintes estatísticas das notas19 : média do desempenho na prova, erro padrão da média, desvio padrão, nota mínima, mediana e nota máxima para cada um de seus componentes.
Tais estatísticas contemplam o total de estudantes concluintes da Área de Tecnologia em Estética e Cosmética inscritos e presentes à prova do Enade/2019, tendo em vista agregações, ou por Grandes Regiões e o país como um todo, ou por Categoria Administrativa, Organização Acadêmica da IES e Modalidade de Ensino.
Em relação aos gráficos de distribuição de notas, o intervalo considerado foi de 10 unidades, aberto à esquerda e fechado à direita, com exceção do primeiro intervalo, [ 0 ; 10 ], fechado em ambos os extremos.
Para os gráficos de distribuição das notas das questões discursivas, foram consideradas mais duas categorias : questão em branco20 e nota zero. 19 Uma definição dessas estatísticas pode ser encontrada no Glossário. 20 Nesse grupo estão incluídas também as respostas classificadas como nulas ou desconsideradas. 6.1 ESTATÍSTICAS BÁSICAS DA PROVA Esta seção apresenta estatísticas selecionadas e histogramas da nota geral ( 6.1.1 ) e de cada componente : Formação Geral ( 6.1.2 ) e Conhecimento Específico ( 6.1.3 ).
São, também, apresentadas estatísticas selecionadas de subpopulações, caracterizadas por Grande Região, Categoria Administrativa, Organização Acadêmica e Modalidade de Ensino. 6.1.1 Estatísticas Básicas Gerais Na Tabela 6.1, são apresentadas as Estatísticas Básicas da prova, por Grande Região, dos estudantes concluintes de Tecnologia em Estética e Cosmética.
A Média das notas da prova, como um todo ( nas seções seguintes serão analisados os Componentes de Formação Geral e de Conhecimento Específico ), foi 45,8, sendo que os estudantes da região Norte obtiveram a Média mais baixa ( 41,5 ), e os da região Sul obtiveram a Média mais alta ( 47,3 ).
As demais Médias foram : 44,8, na região Nordeste, 45,6, na região Sudeste e 45,9, na região Centro-Oeste.
O Desvio padrão para o Brasil, como um todo, foi 11,4, sendo o maior Desvio padrão encontrado na região Norte ( 11,8 ), e o menor, nas regiões Sul e Centro-Oeste ( 11,1 para ambas ), indicando uma dispersão um pouco menor das notas dessa última região.
A região que obteve a maior nota Máxima foi a Sul ( 83,6 ), ao passo que a região que atingiu a menor nota Máxima foi a Centro-Oeste ( 75,9 ).
A Mediana do Brasil, como um todo, foi 45,7, sendo a maior Mediana obtida na região Sul ( 47,4 ), e a menor obtida na região Norte ( 41,4 ).
A nota Mínima do Brasil foi obtida na região Sudeste, onde foi 2,6.
A maior nota Mínima, entre as cinco Grandes Regiões, foi 15,2, nas regiões Norte e Centro-Oeste, as demais regiões obtiveram nota Mínima igual à : Nordeste ( 10,2 ) e Sul ( 9,4 ).
Considerando -se as notas segundo Grande Região, observa -se que existe diferença estatisticamente significativa ao nível de 95 \% entre a Média da região Sul ( maior Média ) e as Médias das regiões Norte, Nordeste e Sudeste.
O comportamento das notas dos estudantes de todo o Brasil pode ser observado no Gráfico 6.1, que apresenta um histograma com a distribuição das mesmas : uma distribuição unimodal com a moda no intervalo ( 40 ; 50 ].
Na Tabela 6.2, são apresentadas informações referentes à Média da nota final, desagregadas por Categoria Administrativa, por Organização Acadêmica e por Modalidade de Ensino.
A Média das notas da prova, como um todo, foi 45,8.
Em relação a Categoria Administrativa, os estudantes das IES Privadas obtiveram Média igual ( 45,8 ) a Média nacional, e os das IES Públicas obtiveram Média inferior a nacional ( 45,6 ).
Observa -se que não existe diferença estatisticamente significativa entre as Médias das notas das IES Públicas e as das Privadas.
A diferença entre as Médias das regiões Sul e Nordeste ( 5,8 ), a maior e a menor Média, é superior à diferença entre a das IES Públicas e Privadas ( 0,2 ), caracterizando -se uma menor diversidade administrativa do que regional.
O Desvio padrão para as IES Públicas ( 11,3 ) foi inferior ao do Brasil, como um todo ( 11,4 ), indicando uma dispersão um pouco menor das notas nesta Categoria Administrativa.
No tocante a Organização Acadêmica, as Universidades e os Centros Universitários obtiveram Média mais alta que a nacional ( 46,2 e 46,5, respectivamente ).
A Média das Faculdades foram menores do que a nacional ( 44,1 ).
Constata -se que existe diferença estatisticamente significativa, ao nível de 95 \%, entre a Média mais alta, 46,5 dos Centros Universitários e a Média mais baixa, 44,1 das Faculdades.
A Média da Modalidade Educação a Distância ( 42,6 ) foi inferior à Média da Educação Presencial ( 46,0 ), e a diferença entre elas é estatisticamente significativa, ao nível de 95 \%. 6.1.2 Estatísticas Básicas no Componente de Formação Geral Na Tabela 6.3, são apresentadas as Estatísticas Básicas em relação ao componente da prova que avalia a Formação Geral dos estudantes concluintes.
A nota de Formação Geral é obtida levando -se em conta os acertos nas sete questões objetivas ( uma foi anulada pela Comissão \textbf{Assessora} de Área - CAA ) e a nota média obtida nas questões discursivas 1 e 2, sendo a nota de cada questão discursiva obtida levando -se em conta 80 \% da nota relativa ao conteúdo e 20 \% da nota relativa ao desempenho linguístico.
Os estudantes de todo o Brasil obtiveram desempenho médio de 31,7.
Quanto à variabilidade, o Desvio padrão das notas dos estudantes do Brasil, como um todo, foi 15,3.
A maior Média foi obtida na região Sul ( 32,8 ), e a menor, na região Norte ( 29,0 ).
As demais Médias foram : 30,9, na região Nordeste, 31,4, na região Sudeste e 31,9, na região Centro-Oeste.
Já o maior Desvio padrão foi obtido na região Nordeste ( 15,8 ) e o menor, nas regiões Norte, Sudeste e Centro-Oeste ( 15,1 ).
Na região Sul o Desvio padrão foi 15,3.
A maior nota no Componente de Formação Geral da prova do Enade/2019 foi 92,6, obtida por, pelo menos, um estudante na região Sul.
A menor nota Máxima foi obtida na região Norte ( 78,8 ).
A Mediana do Brasil, como um todo, foi 30,1, sendo a menor Mediana encontrada na região Norte ( 27,4 ), e a maior, na região Sul ( 31,4 ).
A nota Mínima nessa parte foi zero em todas as regiões.
Considerando -se as notas, segundo Grande Região, observa -se que existe diferença estatisticamente significativa entre Média da região Sul ( maior Média ) e as Médias das notas do Componente de Formação Geral obtidas nas regiões Centro-Oeste, Nordeste e Sudeste.
Observa -se também, que existe diferença estatisticamente significativa entre a Média da região Sul ( maior Média ) e a Média da região Norte ( menor Média ).
No Gráfico 6.2, é possível verificar a avaliação do desempenho dos estudantes do Componente de Formação Geral, a partir do histograma da distribuição das notas correspondentes.
A distribuição é unimodal, com moda em ( 20 ; 30 ].
Nota -se, que, no Gráfico 6.2, as notas apresentam maior dispersão do que as no Gráfico 6.1 ( distribuição das notas da prova ), confirmada pela comparação dos desvios padrões : 11,4 para a nota da prova, como um todo, e 15,3 para o Componente de Formação Geral.
Na Tabela 6.4, são apresentadas as informações referentes ao desempenho dos concluintes do Componente de Formação Geral, em diferentes agregações : Categoria Administrativa, Organização Acadêmica e Modalidade de Ensino.
Observa -se que não existe diferença estatisticamente significativa entre as Médias dos tipos de Categoria Administrativa : a maior Média obtida por estudantes de IES Públicas ( 32,0 ) e a menor, pelos de IES Privadas ( 31,7 ).
Nota -se que, também, há diferença estatisticamente significativa entre a maior das Médias dos tipos de Organização Acadêmica e a menor, o valor maior para os Centros Universitários ( 32,4 ) e o menor para as Faculdades ( 30,5 ).
No que se refere a Modalidade de Ensino, a Educação Presencial apresentou maior Média do que a Educação a Distância ( Médias 31,7 e 31,6, respectivamente ).
Tais Médias não apresentam diferença estatisticamente significativa. 6.1.3 Estatísticas Básicas do Componente de Conhecimento Específico Na Tabela 6.5, são apresentadas as Estatísticas Básicas referentes ao Componente de Conhecimento Específico da área de Tecnologia em Estética e Cosmética.
A nota do componente de Conhecimentos Específicos leva em conta as notas da parte objetiva e a média das notas obtidas nas três questões discursivas da parte discursiva da Área.
A Média do desempenho dos estudantes do Brasil, como um todo, foi 50,5.
A maior Média foi obtida na região Sul ( 52,2 ), e a menor, na região Norte ( 45,6 ).
As demais Médias foram : 50,5, na região Centro-Oeste, 50,3, na região Sudeste e 49,9, na região Nordeste.
Quanto à variabilidade das notas, o Desvio padrão do Brasil, como um todo, foi 12,6, sendo o maior Desvio padrão observado na região Norte ( 13,2 ) e o menor, na região Sul ( 12,2 ).
Os demais desvios foram : 12,4, na região Nordeste, 12,7, na região Sudeste e 12,6, na região Centro-Oeste.
A Mediana das notas dos estudantes de todo o Brasil foi 50,9.
A maior Mediana ocorreu na região Sul ( 52,6 ), e a menor, na região Norte ( 45,7 ).
As demais Medianas foram : 49,7 na região Nordeste e 50,6 nas regiões Sudeste e Centro-Oeste.
A nota Máxima do Brasil, como um todo, foi 90,3, sendo obtida por, pelo menos, um estudante na região Sul.
As demais notas máximas foram : 85,8, na região Norte, 83,9, na região Nordeste, 88,0, na região Sudeste e 85,4 na região Centro-Oeste.
A nota Mínima foi 3,4, sendo obtida por, pelo menos, um estudante na região Sudeste.
As demais notas Mínimas foram : 11,0, na região Norte, 13,6, na região Nordeste, 6,8, na região Sul e 10,2 na região Centro-Oeste.
Observa -se que existe diferença estatisticamente significativa entre a maior Média e as demais, exceto com a Média da região Centro-Oeste.
Assim como nos Gráficos 6.1 e 6.2, no Gráfico 6.3, apresentado a seguir, é possível ser feita uma avaliação do desempenho de concluintes de Tecnologia em Estética e Cosmética, em relação ao Componente de Conhecimento Específico, com um histograma da distribuição das notas correspondentes.
Esta também é uma distribuição unimodal, e o grupo modal é o ( 50 ; 60 ], acima do grupo modal da prova, como um todo, e acima do grupo modal para a Formação Geral.
Na Tabela 6.6, é apresentada uma comparação dos resultados em relação a Categoria Administrativa, a Organização Acadêmica e a Modalidade de Ensino, agora levando -se em conta o desempenho de estudantes do Componente de Conhecimento Específico da prova.
No que se refere a Organização Acadêmica, a maior Média foi a dos Centros Universitários ( 51,2 ), vindo a seguir a das Universidades ( 51,0 ) e, por fim, a das Faculdades ( 48,6 ).
O maior Desvio padrão, e acima do valor para o Brasil, como um todo, foi o dos Centros Universitários e Universidades ( 12,7 ).
As Universidades obtiveram a maior nota Máxima ( 90,3 ).
Os Centros Universitários obtiveram nota Máxima 86,7 e as Faculdades, 88,0.
As Medianas foram : 51,3 nas Universidades, 51,9 nos Centros Universitários e 48,8 nas Faculdades, a menor delas.
A nota Mínima foi 3,4 para Centros Universitários, 6,8 Universidades e 10,2 para as Faculdades.
Observa -se que existe diferença estatisticamente significativa, ao nível de 95 \%, no Componente de Conhecimento Específico entre as notas das Faculdades e das demais Organizações Acadêmicas.
Quanto a Categoria Administrativa, observa -se um comportamento semelhante àquele da parte de Formação Geral e à prova, como um todo, ou seja, não existe diferença estatisticamente significativa entre as Médias das IES Públicas ( 50,1 ) e as das IES Privadas ( 50,5 ).
Neste caso, também, a maior Média foi obtida por estudantes de IES Privadas de ensino.
Quanto ao comportamento das Médias, no que toca a Modalidade de Ensino, há diferença estatisticamente significativa entre a Média da Educação Presencial ( 50,7 ) e a da Educação a Distância ( 46,2 ). 6.2 ANÁLISE DAS QUESTÕES OBJETIVAS Esta seção apresenta estatísticas selecionadas e histogramas das Questões Objetivas de Formação Geral ( 6.2.1 ) e de Conhecimento Específico ( 6.2.2 ).
São, também, apresentadas e comparadas as médias das subpopulações caracterizadas por Grande Região. 6.2.1 Componente de Formação Geral Na Tabela 6.7, são apresentadas as Estatísticas Básicas relativas às sete questões objetivas ( uma questão foi anulada pela Comissão \textbf{Assessora} de Área – CAA ) do componente da prova que abrange a Formação Geral dos estudantes.
A Média do Brasil foi 38,6.
A menor Média foi encontrada na região Norte ( 35,5 ), e a maior, na região Sul ( 39,8 ).
As demais Médias foram : 37,7, na região Nordeste, 38,3, na região Sudeste e 39,3, na região Centro-Oeste.
O Desvio padrão do Brasil foi 19,9, sendo o maior Desvio padrão encontrado nas regiões Nordeste e Centro-Oeste ( 20,4 ), e o menor, nas regiões Norte e Sudeste ( 19,7 ).
O Desvio padrão da região Sul foi 20,0.
As Medianas do Brasil, como um todo, e de todas as regiões foram iguais a 42,9, exceto da região Norte ( 28,6 ).
A nota Máxima 100,0 foi alcançada nas regiões Nordeste, Sudeste e Sul.
As demais regiões tiveram nota Máxima de 85,7.
As notas mínimas ( 0,0 ) foram iguais para todas as regiões.
Na Tabela 6.8, são apresentados o Índice de Facilidade e o Índice de Discriminação ( Ponto-Bisserial ) para cada uma das questões objetivas do Componente de Formação Geral.
Quanto ao Índice de Facilidade, foram usadas as seguintes cores para diferenciar o nível de dificuldade da questão : • Azul para as questões classificadas com índice Muito fácil ( \&gt ; =0,86 ), verde para as questões classificadas com índice Fácil ( 0,61 a 0,85 ), amarelo para as questões classificadas com Médio ( 0,41 a 0,60 ), vermelho para as questões Já quanto ao Índice de Discriminação, foram usadas as seguintes cores para qualificar a questão : • As questões classificadas com índice Fraco receberam a cor vermelha ( \&lt ; =0,19 ), as classificadas com Médio receberam a cor amarela ( 0,20 a 0,29 ), as classificadas com Bom receberam a cor verde ( 0,30 a 0,39 ) e as classificadas com Muito bom ( \&gt ; =0,40 ) receberam a cor azul.
As questões objetivas do Componente de Formação Geral, segundo o Índice de Facilidade, foram assim avaliadas : das sete questões, uma foi anulada, nenhuma teve o Índice de Facilidade classificado como Muito fácil ou Fácil.
Quatro questões foram consideradas com índice de dificuldade Médio, situando -se no intervalo entre 0,41 e 0,53 do Índice de Facilidade, ou seja, houve entre 41,0 \% e 53,0 \% de acertos, enquanto duas questões foram classificadas na categoria Difícil, situando -se no intervalo entre 0,39 e 0,40.
Uma questão apresentou menos de 15 \% de acertos, sendo classificada como Muito difícil.
O Índice de Facilidade variou de 0,11 a 0,53 e o de Discriminação, de 0,22 a 0,49.
Todas as questões, quanto ao Índice de Discriminação, foram tidas como Muito bom, com exceção da questão 5, tendo Índice de Discriminação Médio.
Na Tabela 6.9, é apresentada a distribuição das questões, levando -se em conta, simultaneamente, a classificação dos dois índices.
As seis questões que tiveram Índice de Discriminação Muito bom figuraram entre dois níveis de dificuldade Médio e Difícil : quatro classificadas na categoria Média ( questões 2, 3, 4 e 7 ), do Índice de Facilidade, e duas, na categoria Difícil ( questões 6 e 8 ).
Em particular, a questão 8 foi a que apresentou o maior poder discriminatório, com índice 0,49, e foi considerada Difícil em termos de facilidade, com classificadas com Difícil ( 0,16 a 0,40 ) e roxo para as questões classificadas com Muito difícil ( \&lt ; =0,15 ).
Já quanto ao Índice de Discriminação, foram usadas as seguintes cores para qualificar a questão : • As questões classificadas com índice Fraco receberam a cor vermelha ( \&lt ; =0,19 ), as classificadas com Médio receberam a cor amarela ( 0,20 a 0,29 ), as classificadas com Bom receberam a cor verde ( 0,30 a 0,39 ) e as classificadas com Muito bom ( \&gt ; =0,40 ) receberam a cor azul.
As questões objetivas do Componente de Formação Geral, segundo o Índice de Facilidade, foram assim avaliadas : das sete questões, uma foi anulada, nenhuma teve o Índice de Facilidade classificado como Muito fácil ou Fácil.
Quatro questões foram consideradas com índice de dificuldade Médio, situando -se no intervalo entre 0,41 e 0,53 do Índice de Facilidade, ou seja, houve entre 41,0 \% e 53,0 \% de acertos, enquanto duas questões foram classificadas na categoria Difícil, situando -se no intervalo entre 0,39 e 0,40.
Uma questão apresentou menos de 15 \% de acertos, sendo classificada como Muito difícil.
O Índice de Facilidade variou de 0,11 a 0,53 e o de Discriminação, de 0,22 a 0,49.
Todas as questões, quanto ao Índice de Discriminação, foram tidas como Muito bom, com exceção da questão 5, tendo Índice de Discriminação Médio. uma proporção de 0,40 acertos.
O máximo de acertos foi alcançado pela questão 2, com um Índice de Facilidade de 0,53.
As questões 6 e 8 foram consideradas Difícil, com índices de Facilidade 0,39 e 0,40, respectivamente.
Os índices de Discriminação foram Muito bom para as duas questões ( 0,46 e 0,49 respectivamente ).
No Gráfico 6.4, para exemplificar, analisa -se o comportamento da questão de número 8 de Formação Geral.
Trata -se de uma questão considerada Difícil, em relação à facilidade e a que obteve o maior Índice de Discriminação dessa parte da prova.
Neste gráfico, cada uma das cinco curvas representa o percentual de respostas em determinada alternativa da questão, em função do número de acertos dos estudantes nessa parte da prova ( Formação Geral/Múltipla Escolha ), antes de possíveis eliminações pelo critério Ponto-Bisserial.
A curva em verde corresponde à alternativa B, a alternativa correta para essa questão.
Observa -se que, entre os estudantes com menor número de acertos, nessa parte do exame, a situação mais frequente foi a escolha da alternativa C ( em preto ), incorreta.
Por exemplo, entre os estudantes que acertaram duas questões, 11,9 \% escolheram a alternativa E ( em vermelho ), 16,4 \% escolheram a alternativa D ( em \textbf{laranja} ), 24,9 \% escolheram a alternativa B ( em verde, gabarito ), 10,8 \% escolheram a alternativa A ( em azul ), e 35,6 \%, a C ( em preto ).
Dentre os que acertaram duas respostas entre as questões de múltipla escolha de Formação Geral, 0,4 \% deixou a questão em branco e 0,1 \% marcou mais de uma alternativa, invalidando a questão.
À medida que o número de acertos aumenta, indicando desempenho melhor nessa parte da prova, aumenta, concomitantemente, a proporção de estudantes que selecionaram a alternativa correta B, atingindo 100 \% para os estudantes com sete acertos.
Essa análise permite verificar como a questão discriminou os grupos de desempenho, justificando -se o alto índice ( 0,49 ) obtido na questão.
Cumpre notar que não é possível inferir deste gráfico nem o Índice de Facilidade, que seria uma média da proporção ponderada pela quantidade de estudantes com cada uma das notas ), nem o Índice de Discriminação Ponto-Bisserial, por razão equivalente.
No caso extremo, no qual a grande concentração dos acertos dos estudantes fosse abaixo de quatro, 152 O Índice de Facilidade seria obrigatoriamente abaixo de 20,0 \% ( neste exemplo ).
Caso a concentração fosse em seis acertos ou mais, o índice seria obrigatoriamente acima de 50,0 \%.
Os gráficos relativos às demais questões de Formação Geral constam do Anexo I. 6.2.2 Componente de Conhecimento Específico Na Tabela 6.10, são apresentadas as Estatísticas Básicas em relação às questões objetivas do Componente de Conhecimento Específico da prova de Tecnologia em Estética e Cosmética por Grande Região.
Como será discutido ainda nesta seção, duas questões objetivas do Componente de Conhecimentos Específicos não foram usadas no cômputo das notas por terem sido descartadas pelo critério do Ponto-Bisserial.
Assim, as notas foram calculadas com base em 25 das 27 questões objetivas válidas de Conhecimentos Específicos.
A Média do Brasil deste componente foi 53,8.
A menor Média foi observada na região Norte ( 48,8 ), e a maior, na região Sul ( 55,6 ).
O Desvio padrão de todo o Brasil foi 13,4, sendo o menor Desvio padrão encontrado na região Sul ( 13,1 ), e o maior, na região Norte ( 13,9 ).
A Mediana de todo o Brasil foi 56,0, o mesmo valor da Mediana encontrada nas regiões Sul e Centro-Oeste.
Os valores das demais Medianas foram : Norte ( 48,0 ), Nordeste e Sudeste ( 52,0 ).
A nota Máxima da prova foi 96,0, obtida nas questões objetivas do Componente de Conhecimento Específico, por, pelo menos, um estudante da região Sul.
A nota Mínima foi 4,0, por, pelo menos, um estudante na região Sudeste.
Os valores das demais notas mínimas foram : Norte ( 12,0 ), Nordeste ( 16,0 ), Sul ( 8,0 ) e Centro-Oeste ( 12,0 ).
Na Tabela 6.11, são apresentados os Índices de Facilidade e de Discriminação ( Ponto-Bisserial ) das questões objetivas do Componente de Conhecimento Específico, para os estudantes de Tecnologia em Estética e Cosmética.
Para facilitar a diferenciação das questões, serão usadas as mesmas cores da Tabela 6.8 para as diferentes classificações dos Índices de Facilidade e de Discriminação.
A partir do Índice de Facilidade obtido, pode -se concluir que quase metade das questões objetivas da prova foram considerados, pelo menos, Fácil : das 27 questões válidas, 12 foram classificadas como Fácil ou como Muito Fácil.
Sete questões foram classificadas como Médio, e outras oito consideradas como Difícil.
Nenhuma questão foi classificada como Muito Difícil.
Já quanto aos Índices de Discriminação das questões objetivas do Componente de Conhecimento Específico da prova, obtém -se como resultado a seguinte classificação : 12 das 27 questões foram consideradas como Bom, 13 delas foram classificadas como Médio, e outras duas, como Fraco, sendo 15, por conseguinte, a quantidade de questões nos dois patamares mais baixos de discriminação.
Constata -se, assim, que a prova – no que se refere ao Componente de Conhecimento Específico – possuía média capacidade de discriminar entre aqueles que dominam ou não o conteúdo.
O Índice de Facilidade variou de 0,19 a 0,96 e o de Discriminação, de 0,01 a 0,37.
Na Tabela 6.12, é apresentada a distribuição das questões levando -se em conta, simultaneamente, a classificação dos dois índices.
A questão de número 11 foi a segunda mais difícil entre as 27 questões específicas válidas, com baixo Índice de Facilidade, apenas 22 \% de acertos.
Essa questão apresentou poder discriminatório muito baixo, 0,01 o que comprova ter sido a segunda mais difícil para os estudantes.
Destaca -se, também, a questão 23, com Índice de Facilidade 0,37, o que, em termos percentuais, corresponde a 37 \% de estudantes que responderam acertadamente.
Já 0,01 foi o seu Índice de Discriminação.
Tais questões foram, portanto, pelo critério Ponto-Bisserial, consideradas inadequadas.
Por isso, as questões 11 e 23, foram eliminadas do cômputo da nota final.
A título de exemplo das análises do comportamento das questões objetivas, no Gráfico 6.5, analisa -se a questão 13 do Componente de Conhecimento Específico.
Essa questão foi considerada, pelas respostas dos estudantes avaliados na prova, como uma questão com nível de facilidade Médio, 0,47, ou seja, 47 \% dos estudantes assinalaram, acertadamente, a opção E, correspondente ao gabarito.
Como já comentado, seu Índice de Discriminação foi igual a 0,37, classificado como Bom, o maior valor de discriminação.
Neste gráfico, cada uma das cinco curvas representa o percentual de respostas em determinada alternativa da questão 13, em função do número de acertos dos estudantes nessa parte da prova, antes de possíveis eliminações de questões pelo critério Ponto-Bisserial.
A alternativa correta E, representada no gráfico pela curva em vermelho, foi escolhida em maiores proporções pelos estudantes com desempenho melhor nessa parte da prova.
Já as alternativas incorretas, também denominadas distratores, foram selecionadas, principalmente, por aqueles com notas mais baixas.
No caso de Tecnologia em Estética e Cosmética, como se observa no eixo horizontal do Gráfico 6.5, nenhum estudante acertou mais do que 25 questões dentre as 27 questões objetivas de Conhecimentos Específicos, considerando -se o gabarito original.
Mesmo com a eliminação das duas questões com Índice de Discriminação Fraco do cômputo da nota, nenhum participante atingiu a nota Máxima, igual a 100,0, nessa parte da prova ( visto na Tabela 6.10 ).
Aqueles com nota zero, na sua totalidade, deixaram esta questão em branco ou marcaram mais de uma alternativa, comportamento considerado inválido.
A proporção de estudantes que selecionaram a resposta correta E aumenta gradativamente, chegando a atingir 100 \% para 23 acertos ou mais, enquanto a proporção dos que escolheram alternativas incorretas decai, a partir de três ou quatro acertos, em função do número de acertos nessa parte da prova.
Os gráficos relativos às demais questões do Conhecimento Específico constam do Anexo I. 6.3 ANÁLISE DAS QUESTÕES DISCURSIVAS Esta seção apresenta estatísticas selecionadas e histogramas das Questões Discursivas de Formação Geral ( 6.3.1 ) e Conhecimento Específico ( 6.3.2 ).
São, também, apresentadas e comparadas as médias de subpopulações, caracterizadas por Grande Região.
Cumpre notar que nem todas as questões passam pelo crivo dos corretores.
As respostas depois de digitalizadas passam por um pré-processamento de identificação automático.
Algumas destas questões são definidas como um “ branco automático ” ( que não são enviadas para os professores para correção, a não ser que na segunda fase, a de inspeção visual, seja constatada uma resposta ).
O critério usado neste pré-processamento, é verificar para cada questão, a partir da quantidade de PIXELS, o que provavelmente, é uma resposta em branco.
A partir da separação destes documentos que foram considerados em branco, eles são enviados para uma verificação visual, por uma equipe de colaboradores, devidamente treinados.
Deste modo, quaisquer erros, deste pré-processamento, são encontrados, e as questões preenchidas, são enviadas para a correção dos professores.
Eventualmente, algumas respostas em branco, não são detectadas neste pré-processamento automático.
A presença de pequenos riscos, ou manchas, podem fazer o programa não identificar uma questão como um possível branco.
Estas são enviadas para a correção, mas são classificadas como “ branco ”, pelos professores quando da correção.
No caso da Área de Tecnologia em Estética e Cosmética, foram identificados como brancos “ automáticos ”, 562 respostas de questões discursivas de Conhecimento Específico e 1.853 de Formação Geral.
Pelos professores corretores foram identificadas mais 205 e 251 questões em branco, respectivamente, para Conhecimento Específico e para Formação Geral.
Uma outra situação a ser mencionada é a ocorrências de atendimento a portadores de necessidades especiais.
O anexo X apresenta, para o exame como um todo e para a área de Tecnologia em Estética e Cosmética, por tipo de deficiência, o protocolo usado para permitir a correção das respostas dos alunos em cada situação. 6.3.1 Componente de Formação Geral As análises dos resultados de desempenho dos estudantes de Tecnologia em Estética e Cosmética, nas duas questões discursivas relativas a Formação Geral, encontram -se na Tabela 6.13 e no Gráfico 6.6.
Aqui, analisa -se a média das notas de cada questão discursiva de Formação Geral, que leva em conta as correções realizadas por duas bancas, uma que analisa e pontua aspectos relativos ao conteúdo das respostas, seguindo o padrão de respostas divulgado pelo Inep ( ver Anexo VIII ), outra que avalia o desempenho linguístico expresso pela redação das respostas, seguindo um padrão de respostas específico da área de Língua Portuguesa ( ver Anexo VIII ).
A nota de cada questão é composta de 80 \% da nota da correção de conteúdo e 20 \% da nota de desempenho linguístico.
Na Tabela 6.13, observa -se que a nota Média nesse conjunto de questões foi inferior à obtida nas objetivas.
Os estudantes de todo o Brasil obtiveram, em Formação Geral, Média 38,6, nas questões objetivas e 21,3, nas questões discursivas.
Pode -se notar, também, que o Desvio padrão nesse conjunto de questões foi menor do que o obtido nas objetivas : 19,9, nas questões objetivas e 18,1, nas questões discursivas.
A maior Média foi obtida na região Sul ( 22,4 ), e a menor, na região Norte ( 19,2 ).
A Mediana de todo o Brasil, neste componente, foi 17,0.
Nas regiões Norte, Nordeste, Sudeste e Centro-Oeste, a Mediana foi menor ( 14,5, 16,0, 16,5 e 16,5, respectivamente ), e, na região Sul ( 19,0 ) foi maior que o valor para o Brasil.
A nota Máxima ( 93,5 ) foi obtida na região Sul, sendo 90,0, a Máxima na região Norte ; 81,0, na Nordeste ; 83,5, na Sudeste ; e 83,0, na região Centro-Oeste.
A nota Mínima ( 0,0 ) foi a mesma em todas as regiões do Brasil.
No Gráfico 6.6, está representada a distribuição das notas nas questões discursivas do Componente de Formação Geral.
A moda dessa distribuição ocorre no intervalo [ 0 ; 10 ], com 35,5 \% de frequência, seguida pelo intervalo ( 10 ; 20 ], com 19,4 \%.
Destaca -se que no intervalo [ 0 ; 10 ] se inclui, além da nota zero, a frequência de estudantes que deixaram esse tipo de questão em branco.
Na sequência, os resultados verificados para cada uma das questões discursivas de Formação Geral serão apresentados, estabelecendo -se relações com os conteúdos abordados em cada uma delas.
Os comentários das Bancas de docentes corretores, a respeito do observado na correção das respostas dos estudantes, suas impressões e conclusões serão apresentados junto à análise de cada questão.
Cumpre esclarecer que, tendo em vista que as questões discursivas de Formação Geral são padronizadas, ou seja, constam de todas as provas, os comentários da Banca são os mesmos para todas as carreiras acadêmicas, sendo direcionados a todos os estudantes que participaram do Enade/2019.
A seguir, será analisado o desempenho linguístico dos estudantes da Área de Tecnologia em Estética e Cosmética nas duas questões discursivas de Formação Geral do Enade/2019, apresentando -se os resultados obtidos e os comentários da banca de correção para cada questão. 6.3.1.1 Análise de Conteúdo da Questão Discursiva 1 do Componente de Formação Geral Os dados de Tecnologia em Estética e Cosmética, obtidos a partir das respostas à questão 1, encontram -se na Tabela 6.14 e no Gráfico 6.7.
Nessa questão – de desempenho muito semelhante ao verificado na outra questão de Formação Geral – os estudantes de todo o Brasil, tiveram Média 21,2.
A maior Média para a questão 1 foi obtida na região Sul ( 23,6 ), e a menor, na região Norte ( 16,8 ).
Quanto à variabilidade das notas, o Desvio padrão de todo o Brasil foi 29,5.
O menor Desvio padrão foi obtido na região Norte ( 27,1 ), e o maior Desvio padrão foi obtido pela região Sul ( 31,2 ).
A Mediana do Brasil, como um todo, foi zero, sendo igual em todas as regiões.
As notas Máxima e Mínima da questão discursiva 1 foram as mesmas para todas as regiões do Brasil, respectivamente, 100,0 e zero, com exceção das regiões Norte e Nordeste que tiveram nota Máxima de 95,0.
No Gráfico 6.7, mostra -se a distribuição das notas na questão discursiva 1 do Componente de Formação Geral.
Observa -se que 14,0 \% dos participantes deixaram em branco a resposta a essa questão.
Somados à frequência dos que, tendo respondido, receberam nota zero, chega a 57,6 \% o contingente daqueles que zeraram a questão.
O intervalo ( 40 ; 50 ] se caracteriza como um máximo local com frequência 10,9 \%, seguido de perto do intervalo ( 20 ; 30 ], com 9,9 \%. 6.3.1.2 Comentários sobre a correção de Conteúdo das respostas à Questão Discursiva 1 No enunciado da Questão 1, que era constituída por dois textos e um comando, abordavam -se dois itens da lista de temas de referência estabelecidos pelo Inep, pela Portaria nº 510, de 31 de maio de 2019, que dispõe sobre o componente de Formação Geral : ( VIII ) meio ambiente, sustentabilidade e intervenção humana e ( IX ) cidades, habitação e qualidade de vida.
No primeiro texto, apresentavam -se estatísticas sobre desastres naturais ocorridos no Brasil e a eles associavam -se eventos recorrentes, tais como estiagens, secas, inundações \textbf{bruscas} e alagamentos.
O segundo texto complementava o primeiro, nele constava um relatório da ONU no qual é identificada a necessidade de se reduzirem os níveis existentes de riscos que favorecem os desastres, apresentando o fortalecimento da resiliência social, ambiental e econômica como uma das soluções para que as cidades consigam conviver com esses fenômenos naturais.
No comando, por outro lado, não se solicitava reflexão sobre os riscos de futuros desastres, temática do segundo texto.
Solicitava -se reflexão sobre ações após a ocorrência de um desastre : “ duas propostas de intervenção no âmbito da sustentabilidade socioambiental de modo a contemplar ações de recuperação ou de restauração após a ocorrência de desastres ”.
Assim, apesar de os textos de contextualização serem claros e conterem informações relevantes – fatos de conhecimento público e de grande divulgação –, havia uma diferença de abordagem entre os dois e o comando em relação à linha do tempo da ocorrência de desastres.
Esse fato aumentou o grau de dificuldade da questão.
Entretanto, os conhecimentos exigidos e o nível de profundidade esperado pelo padrão de resposta eram compatíveis com a formação universitária.
Ao se exigirem propostas de intervenção envolvendo a sustentabilidade socioambiental, a questão possibilitou a avaliação, em primeiro lugar, da capacidade de elaborar e apresentar propostas de ação e de intervenção, a partir da análise de um contexto, e, em segundo lugar, da capacidade de propor soluções viáveis e inovadoras na resolução de uma situação-problema.
Na solicitação do comando estava implícita a exigência de redação de um texto em que se apresentasse cada proposta, o que deveria incluir informações suficientes para sua compreensão, ou seja, como as ações seriam \textbf{concretizadas}, por quem seriam realizadas ou que benefícios trariam à população ou à área atingida.
Na questão, exigia -se, portanto, a capacidade de o estudante expor e argumentar, de modo a defender ou justificar suas propostas, de forma coerente e articulada.
Além disso, pelo conteúdo da resposta era possível avaliar o estágio de comprometimento do estudante com as questões sociais e ambientais.
A dicotomia entre o antes e o após desastre, criada no confronto entre os textos e o comando da questão, permitiu aquilatar também a aptidão do estudante para ler e interpretar textos não lineares, em que são exigidas atenção e percepção de diferenças sutis.
No entanto, tal dicotomia gerou respostas que não estavam em conformidade com as linhas de ação apontadas no padrão de resposta, respostas erradas por terem foco em ações de prevenção.
O padrão de respostas foi apresentado em linguagem simples, compatível com o nível dos formandos.
O padrão tinha flexibilidade e extensão, visto que apresentava inúmeras ações de recuperação e de restauração, no âmbito da sustentabilidade socioambiental, sem exauri -las, como passíveis de abordagem pelos estudantes em suas respostas.
Também não havia exigência de conhecimentos específicos sobre qualquer matéria, ficando restrito aos limites de conceitos básicos sobre meio ambiente, sustentabilidade e intervenção humana em emergências provocadas por desastres, além de atualização com assuntos amplamente tratados na mídia nacional.
O padrão cobriu diversas áreas de atividades e categorizou -as em ações psicossociais, econômicas e sociais, ambientais, de infraestrutura e sistêmicas, com destaque para as ocorrências de estiagens e secas, porque são fenômenos praticamente permanentes em diversas regiões do país.
Esse universo de intervenções criou amplo espaço para o enquadramento das respostas dadas pelos estudantes, incluindo enfoques variados do ponto de vista das diferentes formações e especificidades das áreas dos participantes do Enade 2020.
O conteúdo das respostas analisadas na correção foi bem coberto pelas propostas sugeridas no padrão.
Da mesma forma, as respostas corretas cobriram praticamente todo o espectro levantado pelo padrão.
Na maior parte das respostas que não atenderam ao padrão, propunham -se ações de prevenção, ou seja, procedimentos anteriores aos desastres e, portanto, equivocadas.
Nas respostas alinhadas com o padrão, preferencialmente, havia sugestões de ações, dentre as quais se destacaram : recuperação da flora nativa, controle da qualidade da água, organização e treinamento de grupos de voluntários para ações de prevenção de surtos e epidemias, orientação para evitar novos desabamentos e redução de efeitos dos alagamentos, concessão de benefícios como aluguel social e financiamento para recuperação de residências ou reativação de negócios, projetos de realocação de desabrigados e suporte psicológico, médico e social às pessoas e comunidades afetadas pelas calamidades.
Nas respostas envolvendo situações de estiagem e seca, apresentavam -se ações que tanto poderiam ser consideradas de prevenção quanto de recuperação, pelas características sistêmicas desses desastres.
Por isso, propostas como distribuição de água por caminhões pipa ou perfuração de poços foram consideradas corretas.
Apesar de fora do contexto da questão, foram bastante citadas sugestões de penalização dos supostos culpados pelo desastre com multas, bem como recompensa às pessoas ou empresas que prestassem socorro nessas situações.
Por outro lado, apesar de presente no padrão, ações de resgate do patrimônio histórico e cultural praticamente não foram citadas como resposta.
Em um conjunto grande de respostas, os estudantes se limitaram a apresentar apenas uma proposta e, nesse caso, preferencialmente, de ajuda a desabrigados ou de recuperação de matas ou de ações vinculadas à seca e à estiagem.
Dentre as respostas consideradas corretas, muitas eram pobres na argumentação e, por vezes, confusas, apesar de estarem em conformidade com o padrão.
Foi percebida certa dificuldade para se organizar a argumentação no contexto do comando.
Em muitas respostas, citavam -se ações isoladas sem que se caracterizasse o seu desdobramento.
Como já comentado, grande parte dos estudantes, em desacordo com o comando da questão, fixou -se em ações de prevenção, anteriores aos desastres.
Houve respostas de excelente conteúdo, mas totalmente enquadradas no contexto preventivo.
Esse viés de prevenção também incluiu respostas fracas em que se parafraseavam os textos do enunciado, mencionando -se de forma superficial inundação, alagamentos, seca e redução dos níveis de risco.
A interpretação incorreta pode ter sido induzida pelo enunciado, mas o principal indício parece ser de falta de atenção na leitura do comando.
As maiores fragilidades na formação dos estudantes, apontadas pelos corretores, foram a leitura e interpretação da questão e a capacidade de expressão linguística.
As respostas com geração de ideias próprias ocorreram em escala bem menor do que aquelas em que os estudantes usaram do recurso da paráfrase pobre sobre o texto do enunciado, revelando dificuldade para criar e articular ideias a respeito dos temas da questão.
Além disso, pelas sugestões que envolviam aspectos econômicos e financeiros revelou -se um nível razoável de desconhecimento das atribuições e responsabilidades das diversas instâncias do governo e de como são encaminhados os recursos necessários ao socorro de populações atingidas por desastres.
Como aspectos relevantes ainda não considerados, podem ser citados o baixo percentual de protestos e a falta de viés crítico nas respostas, talvez pelo fato de o comando ser bastante objetivo. 6.3.1.3 Análise de Conteúdo da Questão Discursiva 2 do Componente de Formação Geral A Tabela 6.15 mostra que o desempenho médio dos estudantes, na questão discursiva 2 ( média 11,2 ), foi inferior ao obtido na questão discursiva 1 ( média 21,2 ).
A região Norte foi aquela cuja Média foi maior ( 12,1 ), a de menor Média foi a região Centro-Oeste ( 9,7 ).
Quanto à variabilidade das notas, o Desvio padrão de todo o Brasil foi 20,4, inferior ao obtido na questão discursiva 1 ( 29,5 ).
O maior desvio nessa questão foi obtido na região Norte ( 21,7 ), enquanto o menor foi obtido na região Sul ( 18,5 ).
A Mediana de todo o Brasil foi 0,0, a mesma todas as regiões.
As notas Máxima ( 100,0 ) e Mínima ( 0,0 ) foram as mesmas em todas as regiões, com exceção das regiões Sudeste e Centro-Oeste que obtiveram nota Máxima de 75,0.
No Gráfico 6.8, mostra -se a distribuição das notas na questão discursiva 2 do Componente de Formação Geral.
Observa -se que 25,5 \% dos estudantes deixaram a resposta à questão em branco.
Chega a 17,9 \%, se agregado à frequência daqueles que, tendo respondido, receberam nota zero.
Observa -se um máximo local no intervalo ( 30 ; 40 ], com 10,0 \% do total de participantes.
Nota -se, ainda, que as notas ficaram levemente menos dispersas em comparação às da questão discursiva de número 1, o que pode ser constatado, também, pela comparação do desvio padrão das notas da questão discursiva 2 ( 20,4 ) e o das notas da questão discursiva 1 ( 29,5 ). 6.3.1.4 Comentários sobre a correção de Conteúdo das respostas à Questão Discursiva 2 Na Questão 2, o tema era a identificação de impactos científicos e econômicos do desenvolvimento do Brasil na área de Ciência e Tecnologia ( C\&T ), bem como o papel dos principais atores responsáveis por impulsionar esse desenvolvimento.
O tema escolhido era relevante por contemplar os referenciais divulgados pelo Inep para esse exame : Ciência, tecnologia e inovação – Estado, sociedade e trabalho, assuntos supostamente acessíveis aos formandos do nível universitário brasileiro de qualquer área.
A análise da questão demonstra sua conformidade com as diretrizes de avaliação de competências e habilidades dos estudantes.
A explicitação de ganhos vinculados a dois campos diferentes, no contexto de C\&T, e como consequência de colaboração interinstitucional, proporciona a verificação da capacidade não só de ler e interpretar textos, como de narrar, de forma consistente e articulada, o resultado de reflexão sobre necessidades em contextos diversos.
Além disso, as descrições de articulações e escolhas de ganhos que se \textbf{julgassem} relevantes permitiram que se avaliasse o nível de maturidade do conhecimento sobre os sistemas científicos, tecnológicos e econômicos brasileiros.
O enunciado era claro, e sua linguagem, compatível com o que se espera de formandos de cursos superiores.
Era constituído por uma figura e um texto de contextualização que auxiliavam na compreensão do que foi solicitado pelo comando.
No texto de apresentação, comparava -se a posição do Brasil com a de outros países em indicadores de produção e utilização de conhecimento e de novas tecnologias.
Na figura ilustrava -se a relação entre as principais entidades presentes no contexto de C\&T – as instituições de ensino superior, as empresas e o setor público.
Havia o pressuposto, portanto, de que os egressos fossem capazes de sustentar ideias articuladas e níveis de argumentação, a partir dos elementos textuais fornecidos pela questão, a saber, a figura e o texto, elementos de amparo à elaboração da resposta.
No comando, interligando -se texto e figura, pedia -se para o estudante citar dois ganhos possíveis para o campo científico e dois ganhos possíveis para o campo econômico, resultantes de uma boa articulação entre os entes representados na figura.
Aparentemente, a tarefa solicitada seria simples.
Entretanto, exigia reflexão do estudante em três aspectos : ( i ) a exigência de um processo de articulação entre, pelo menos, duas das entidades explicitadas na figura ; ( ii ) a criação de duas propostas de ganhos em dois setores diferentes, o científico e o econômico, ambos no contexto de C\&T ; e ( iii ) que os ganhos fossem consequência da colaboração de, pelo menos, dois dos entes.
Essa singularidade do enunciado, elevou a dificuldade da questão.
O comando permitia duas possibilidades de interpretação da articulação entre os entes para a estruturação das respostas.
Uma seria o estudante interpretar que seria necessário explicitar em sua resposta o inter-relacionamento gerador dos ganhos ; a outra, induzida pela figura, seria considerar que a articulação entre entidades já estava dada pelo enunciado, deixando ao estudante apenas o registro dos ganhos.
É presumível que a intenção inserida no comando, e refletida no padrão de resposta, fosse a primeira opção, mas a segunda interpretação era plausível.
O padrão corroborava a primeira opção, pois parecia esperar dos formandos a identificação de uma relação de \textbf{causa-efeito} da ação conjunta de dois ou mais entes e do relato das consequências científicas e econômicas dessa ação conjunta.
O padrão de resposta era abrangente e redigido em linguagem simples, apropriada, com nível de profundidade coerente com a formação geral esperada dos respondentes.
Foi apresentado na forma de duas listas extensas, com três temas científicos e oito econômicos, expressos de forma genérica, cada tema completado com exemplos de ganhos passíveis de citação pelos estudantes.
A variedade de temas abordados no padrão proporcionou a flexibilidade e a extensão necessárias para orientar critérios de avaliação.
As respostas mantiveram -se, em sua grande maioria, no escopo antecipado pelo padrão de respostas e no universo de linguagem adequado ao cenário de C\&T da questão.
Nesse contexto da linguagem, foi peculiar constatar a existência de grupos de respostas semelhantes e reveladoras da área de origem dos estudantes.
Por exemplo, desenvolvimento de vacinas, pele de tilápia no tratamento de queimaduras e criação de projetos de máquinas industriais de maior produtividade.
A redação, entretanto, foi quase sempre mais simples do que o esperado pelo padrão de resposta, sem desenvolvimento de argumentação mais aprofundada e, às vezes, sem qualquer tipo de argumentação.
Há dois motivos para esse discurso mais direto dos estudantes.
O mais importante veio do próprio comando da questão, em que constava o verbo citar ; liberando a resposta de um compromisso maior com a dissertação.
O segundo, talvez induzido pela figura, veio da compreensão de que a articulação já estaria pressuposta e que bastava a apresentação dos ganhos.
As respostas mais bem avaliadas mantiveram -se no contexto do padrão, mas seguiram duas vertentes ; nelas, os estudantes consideravam ou não a articulação como pré-estabelecida.
Aquelas em que os estudantes fizeram a associação no próprio texto da resposta tinham um conteúdo mais rico e mais bem desenvolvido.
Outra característica das respostas corretas, na conformidade com o padrão de resposta, foi a separação precisa entre os ganhos científicos e econômicos, mantendo -se o contexto de C\&T do texto de apresentação.
Nas respostas equivocadas, seguindo essas duas linhas de interpretação, os estudantes não fugiram dos temas existentes no padrão de resposta, mas não os exploraram adequadamente.
Foram vários os motivos.
Grande parte formou um conjunto de respostas muito sucintas, que não chegavam a qualificar, com clareza, os ganhos como científicos ou econômicos.
Outro conjunto misturou conceitos de educação com pesquisa, ciência com produção, economia com emprego ou trabalho, dificultando a caracterização do ganho e afastando a resposta do padrão.
As propostas de ganhos com viés econômico geraram maior quantidade de não conformidades do que aquelas que trataram do viés científico.
Isso foi resultado de uma dificuldade demostrada pelos estudantes em evidenciar o vínculo do ganho econômico com o contexto de C\&T, quase sempre apresentado com exemplos de vantagem obtida ou concedida pelo setor público a representantes das outras entidades apontadas no enunciado, tais como : aumento de emprego, mais vagas nas universidades ou redução de impostos.
A análise das respostas revelou que, na maioria delas, os ganhos eram descritos de forma genérica, sem qualquer articulação ou exposição de motivos, sem indicar como ou de de onde surgiam, sem especificar o contexto do assunto ou da área de conhecimento envolvida.
Na maioria dos casos, a ideia de citação levou à não redação de frases, mas à menção de palavras ou expressões.
Exemplos frequentes foram : aumento de pesquisa, melhoria de conteúdo, melhoria de estrutura.
Mesmo relevando a interpretação de que essa colaboração já vinha do enunciado, foi surpreendente que em grande parte dessas respostas não se descrevesse como alguma cooperação seria necessária ou deveria ter sido estabelecida para a geração desses ganhos.
Considerar a articulação vinda do enunciado não dispensava essa revelação.
Por outro lado, houve respostas em que estudantes discutiram a importância da articulação, mas, curiosamente, sem citar os ganhos gerados por elas.
Essa dúvida sobre a articulação ficou bem caracterizada nas respostas em que o ganho descrito não vinha de uma parceria, mas simplesmente da ação de uma das entidades em benefício de outra.
No conjunto das respostas sobre ganhos no campo científico, houve prevalência de citações sobre a necessidade de investimentos em pesquisa.
Essa é uma demanda real, bastante discutida no meio acadêmico.
Assim, respostas que apresentaram uma estrutura do tipo “ maiores investimentos/financiamentos em pesquisas ” e valorização dos pesquisadores foram muito comuns.
Geralmente, elas se complementavam com “ pesquisas realizadas pelas IES ”.
Todavia, a argumentação inconsistente dos estudantes evidenciou que existe dificuldade não só em distinguir as funções de pesquisa das de educação, como também a respeito de como ocorrem os processos científicos e tecnológicos.
Poucos demonstraram compreender que os resultados obtidos na área de C\&T são incorporados pelos diversos campos do saber ao longo do tempo e que se consolidam nas técnicas e procedimentos apresentados nos cursos de graduação.
É \textbf{compreensível} que grande parte dos egressos não tenha se envolvido e não venha a seguir o caminho da ciência ; além disso, muitos deles fizeram cursos em EAD ou cursos de graduação tecnológica, não tendo a chance de conviver com esse contexto de C\&T.
No entanto, pelas próprias diretrizes do exame, supõe -se que a formação universitária deveria contemplar a compreensão dessas atividades.
Em uma visão incompleta da ciência, dois fatos importantes ainda foram observados, ambos equivocados : o entendimento tácito de que o desenvolvimento da pesquisa científica é de responsabilidade única das IES e o vínculo desse desenvolvimento a bolsas em cursos de mestrado e doutorado e até mesmo de iniciação científica.
Nessa linha, em quantidade significativa de respostas exemplificou -se o investimento na educação como ganho científico.
No que tange ao item sobre os aspectos econômicos, em boa parte das respostas, os estudantes não estabeleceram relação com o contexto científico demandado pelos elementos da questão.
Nas respostas, não se associavam os ganhos econômicos a qualquer processo científico ou tecnológico.
Muitas vezes, refraseavam a ideia já apresentada como resposta para o campo científico.
Foi majoritário o raciocínio simplista de que ganho econômico é igual a lucro.
Dessa forma, chegaram a citar a “ adoção de mão de obra de estagiários, por elas serem mais baratas que os empregados formais ” como ganho econômico.
Essa concepção de lucro apareceu bastante na proposta de produção e venda de patentes, inclusive para o exterior, talvez porque a palavra patente tenha sido mencionada no texto do enunciado.
Apresentado como resultado das três combinações possíveis de parceria, muitas vezes, esse ganho com patente era usado, na mesma resposta, como ganho científico.
Curiosamente, a desburocratização da sua concessão também foi apresentada como ganho econômico.
Outro ganho econômico bastante citado, por melhorar a lucratividade das empresas, foi a capacitação profissional realizada nas IES, proposta em que há um \textbf{beneficiário} e um doador, mas não há a parceria exigida pelo comando, nem o vínculo com o campo científico tecnológico.
Além disso, as sugestões que envolviam os ganhos econômicos revelaram desconhecimento das atribuições do governo e de como são encaminhados os recursos necessários ao estímulo às atividades cientificas ou econômicas.
A ausência de citações envolvendo órgãos de fomento foi reveladora desse aspecto.
As respostas dos estudantes e o desempenho global retratado na concentração em notas baixas revelam que a formação de nível superior do Brasil, ao menos nos cursos envolvidos nesta edição do Enade, não tem garantido a compreensão do papel da ciência, do papel dos atores que a promovem e de suas consequências socioeconômicas.
Observou -se um desconhecimento generalizado do processo como um todo.
Essa lacuna prejudica tanto o processo de renovação, como o entendimento de que a pesquisa é atividade que requer investimento alto e contínuo, e cujo tempo de retorno destoa do que é praticado no mercado corporativo.
A falta de compreensão do sistema de C\&T, mesmo em nível geral, levou a um número expressivo de respostas em que se indicaram apenas relações de causa e efeito de nível extremamente superficial, do tipo “ a relação entre A e B impacta positivamente a produção científica aumentando as pesquisas ”.
Tal resposta não demonstra compreensão sobre o papel de cada ente, nem de como a relação entre eles se desenvolve para impactar a produção científica. 6.3.1.5 Considerações Finais sobre a correção de Conteúdo das Questões Discursivas do Componente de Formação Geral O principal aspecto detectado nas duas questões foi a dificuldade de os estudantes interpretarem corretamente o que estava sendo solicitado no comando da questão, causada ou por desatenção ou por desconhecimento mais profundo do contexto.
Esse fato foi o principal responsável pelo resultado aquém do esperado para o nível \textbf{formativo} em termos de pontuação.
Os textos daqueles que responderam em conformidade com o padrão demonstraram um universo de argumentação muito limitado, pobre em vocabulário.
O recurso mais comum foi buscar nos textos introdutórios os elementos para a construção de suas frases, o que criou uma peculiaridade interessante, com grupos de respostas muito parecidas.
Encadear e relacionar argumentos também foi uma deficiência identificada pelos docentes da banca de correção.
Os argumentos, quando apresentados, eram redigidos de forma confusa e sem constituir unidade.
Os estudantes apontaram como viés a preocupação social sobre os problemas abordados ; são bons exemplos as propostas de amparo psicológico e material para as vítimas de desastres e o desenvolvimento de vacinas.
Por outro lado, apontaram deficiências consideráveis na compreensão dos papéis que o setor público, as empresas e as IES desempenham na sociedade e no conhecimento tanto do sistema brasileiro de C\&T como da distinção e da articulação entre ensino, pesquisa e extensão. 6.3.1.6 Análise de Língua Portuguesa da Questão Discursiva 1 do Componente de Formação Geral Os dados de Tecnologia em Estética e Cosmética, obtidos a partir das respostas à questão discursiva 1 do Componente de Formação Geral, no que tange à Língua Portuguesa, encontram -se na Tabela 6.16 e no Gráfico 6.9.
Nesse aspecto, os estudantes de todo o Brasil obtiveram Média 46,8.
A maior Média em relação à Língua Portuguesa foi obtida na região Centro-Oeste ( 49,6 ), e a menor, na região Norte ( 44,6 ).
Quanto à variabilidade das notas, o Desvio padrão de todo o Brasil foi 23,9.
O menor Desvio padrão foi obtido na região Sul ( 23,0 ) e o maior Desvio padrão foi obtido na região Norte ( 26,8 ).
A Mediana das notas de Língua Portuguesa na questão 1 foi 55,0 para o Brasil, como um todo, a mesma obtida nas regiões Nordeste, Sudeste e Centro-Oeste.
A Mediana das regiões Norte e Sul foi 50,0.
A nota Máxima para todo o Brasil foi 95,0, com, pelo menos, um estudante tirando essa nota nas regiões Sudeste e Sul.
Nas regiões Norte, Nordeste e Centro-Oeste a nota Máxima ( 90,0 ).
Já a nota Mínima foi zero em todas as regiões do país.
No Gráfico 6.9, mostra -se a distribuição das notas de Língua Portuguesa do Componente de Formação Geral da questão 1.
Observa -se que a maior frequência, igual a 23,9 \%, corresponde à dos estudantes que obtiveram nota no intervalo ( 10 ; 20 ], o intervalo modal da distribuição, seguido do intervalo [ 0 ; 10 ], com frequência 19,3 \%.
Destacam -se, também, os estudantes que deixaram a questão em branco, representando 14,0 \% do total. 6.3.1.7 Comentários sobre a correção de Língua Portuguesa das respostas à Questão Discursiva 1 A Questão 1 se contextualizava em torno do tema"desastres ambientais", tomando como base dois fragmentos de textos acadêmicos publicados em sites oficiais de instituições de ensino superior na internet.
O texto motivador 1 cita um levantamento da ocorrência de desastres naturais associados a fenômenos naturais no Brasil e o texto motivador 2 cita um relatório da ONU sobre a necessidade de minimizar os riscos e os impactos dos futuros desastres naturais.
Com base nessa reflexão inicial, o enunciado solicitava que o participante apresentasse duas propostas de intervenção no âmbito da sustentabilidade socioambiental, de modo a contemplar ações de restauração ou recuperação após a ocorrência desses desastres.
Apesar de não explicitar a exigência de elaboração de um texto dissertativo, o enunciado detalhava várias informações que deveriam compor as duas propostas a serem apresentadas na resposta.
Para atender a essas solicitações, o participante deveria, para cada proposta, elaborar um texto expositivo para detalhar o tipo de desastre ambiental a que estava se referindo, justificar a proposta de intervenção em função do objetivo de garantir a sustentabilidade ambiental e descrever as ações de restauração ou recuperação inseridas na proposta.
Para atender a essas exigências, o vocabulário utilizado deveria ter algum perfil mais técnico e manter coerência com a situação envolvida.
A análise do desempenho linguístico esperado de um estudante que está concluindo seu curso superior apoia -se no princípio de que ele deve utilizar seus conhecimentos sobre o tema para estruturar seus textos de acordo com as características formais da situação comunicativa.
Essa configuração determina exigências quanto à adequação da seleção vocabular, ao desenvolvimento coerente do conteúdo, à estruturação sintática dos períodos, à organização lógica das ideias, à utilização de procedimentos de encadeamento textual e de referenciação, à obediência às exigências morfossintáticas próprias da modalidade escrita da norma-padrão, ao respeito às regras ortográficas e às regras de acentuação gráfica.
Assim, tomando como parâmetro o que pode ser considerado como um desempenho linguístico ideal, a avaliação das respostas procurou mapear detalhadamente o domínio dos participantes quanto aos recursos disponíveis na Língua Portuguesa em situações de escrita formal.
Dessa forma, levaram -se em consideração, no padrão de resposta, as seguintes competências linguísticas próprias da escrita formal culta, responsáveis pelas condições de textualidade : domínio das convenções ortográficas, domínio dos procedimentos de estruturação textual, domínio das convenções de caráter morfossintático e domínio da seleção vocabular.
Essas quatro competências foram reunidas em três itens de avaliação, atendendo aos aspectos explicitados anteriormente : - domínio das convenções ortográficas : grafia de vogais e consoantes, uso de maiúsculas e minúsculas, emprego do hífen e acentuação gráfica ; - domínio dos procedimentos de estruturação textual do ponto de vista microestrutural : organização interna dos períodos, emprego de conectores para a articulação lógica entre os períodos e entre os parágrafos ; emprego de marcas de referenciação lexical e pronominal ; utilização dos sinais de pontuação que contribuem para a organização lógica da frase ; paragrafação do texto ; - domínio das convenções de caráter morfossintático estabelecidas como modelares do ponto de vista da modalidade escrita formal culta da Língua Portuguesa : concordância nominal e verbal, regência nominal e verbal, colocação pronominal, flexão nominal e verbal, correlação entre tempos e modos verbais, ausência de marcas de oralidade.
O domínio da seleção vocabular foi incorporado a essa última competência, tendo em vista a intersecção entre as duas do ponto de vista das exigências do registro formal da modalidade escrita culta.
Os critérios para análise de cada uma das competências que envolvem o domínio das convenções ortográficas, da estruturação textual, das convenções morfossintáticas e da seleção vocabular da escrita formal culta da Língua Portuguesa vêm explicitados a seguir.
Espera -se que o participante :  utilize repertório vocabular diversificado, sem repetição inadequada de termos ;  recorra a vocabulário preciso em relação ao tema da questão ;  complemente o sentido das palavras por meio de expressões modificadoras ( exemplo : “ Desenvolvimento de projetos ” – sentido incompleto no que diz respeito à palavra “ desenvolvimento ” e à palavra “ projetos ” ) ;  empregue um vocabulário mais formal, com ausência de marcas da oralidade, como termos de sentido muito genérico ( “ coisa ”, “ negócio ”, “ você ” ) e termos de registros mais informais ( como gírias, jargões, frases feitas, ditados populares, termos regionais ) ;  seja coerente na combinação vocabular.
O padrão de resposta foi aplicado integralmente para a avaliação das respostas da Questão 1 ( aspectos ortográficos, textuais, morfossintáticos e vocabulares ).
Apesar de o enunciado não fazer referência à exigência de elaboração de um texto dissertativo sobre o tema, a grande maioria dos participantes procurou desenvolver as duas propostas de intervenção de modo detalhado, o que permitiu a avaliação dos quatro aspectos linguísticos do padrão de resposta.
Observou -se, entretanto, uma significativa diversidade de situações de resposta desde a citação de duas propostas de ação sob a forma de frases nominais ( exemplo :"esclarecimento à população sobre os riscos ambientais") ou verbais ( exemplo :"evitar a propagação do efeito estufa"), sem desenvolvimento ou justificativas, até a construção de respostas com detalhamento das propostas apresentadas, acompanhadas por justificativas.
Com base na avaliação das respostas da Questão 1, pode -se considerar que o desempenho dos participantes foi de nível mediano, em sua maior parte, com desvios pontuais nas quatro competências analisadas.
As respostas que foram avaliadas como fracas são de dois tipos principais : respostas itemizadas, com uma frase curta para cada proposta apresentada ; ou respostas constituídas por pequenos textos, com muitos desvios nas quatro competências, principalmente desestruturação sintática e falta de domínio das convenções morfossintáticas.
Nas respostas que foram avaliadas como medianas, os estudantes apresentaram uma estruturação textual relativamente adequada, com ocorrências pontuais de desvios nas quatro competências.
Entretanto, em todos os aspectos, também houve estruturas em que se respeitaram as características da modalidade escrita formal do padrão culto da Língua Portuguesa.
Nas respostas que foram avaliadas como boas, os estudantes apresentaram desenvolvimento textual fluente e coerente, vocabulário formal e técnico, domínio das convenções morfossintáticas e ortográficas.
Na sequência, para cada uma das competências analisadas durante a correção das respostas dos participantes para a Questão 1, destacam -se as principais características observadas.
Aspectos ortográficos – O desempenho dos participantes revelou uma diferença muito grande nos dois aspectos analisados, como tem sido observado em edições anteriores do Enade : baixo índice de desvios da grafia padrão e grande índice de desvios de acentuação.
Em vários casos, ocorre ausência quase completa de acentuação gráfica.
Com base na análise das respostas, conclui -se, portanto, que, de uma maneira geral, os participantes dominam as convenções relativas à grafia das palavras, mas desconhecem, ou desprezam, as convenções relativas à acentuação gráfica.
Seguem exemplos de algumas ocorrências : a ) eliminação do acento indicador da sílaba tônica : - palavras proparoxítonas : “ tecnica ” ( por"técnica"), “ proposito ” ( por propósito"), “ catastrofe ” ( por \textbf{catástrofe}"), “ ambito ” ( por"âmbito"),"sustentavel"( por"sustentável") ; - palavras paroxítonas terminadas em ditongo crescente : “ reincidencia ” ( por"reincidência"),"industrias"( por"indústrias") ; - palavras oxítonas : “ comite ” ( por"comitê") ; b ) uso indevido do acento gráfico, \textbf{localizado} em sílaba átona : “ ecônomia ” ( por"economia"), “ prevênção ” ( por"prevenção"), “ intervênção ” ( por intervenção"),"reciclavél"( por"reciclável") ; c ) omissão da cedilha, do til e do traço do T :"seguranca"( por"segurança"),"realizacao"( por"realização"),"proibicao"( por"proibição").
Há, inclusive, respostas com ausência total de cedilha ou til ; d ) desvios de grafia :"concientizar"( por"conscientizar"), “ atravéz ” ( por"através"), “ intrui ” ( por “ instrui ” ), “ intruídos ” ( por “ instruídos ” ), “ extrangeiro ” ( por"estrangeiro"), “ envestir ” ( por"investir"), “ tecnoligia ” ( por"tecnologia"), “ evulução ” ( por"evolução"), “ consiliar ” ( por"conciliar"), “ extrutura ” ( por"estrutura") ;"serto ” ( por"certo"),"siguinificativos ” ( por"significativos"), “ esportamos ” ( por"exportamos") ; e ) omissão de sílaba : “ natuza ” por “ natureza ” ; f ) desvios de segmentação : “ apartir ” ( por"a partir"), “ a cerca ” ( por “ acerca ” ), “ oque ” ( por"o que") ; g ) uso indevido de inicial maiúscula :"Boletos"; h ) ao contrário do que se esperava, apareceram raramente as abreviaturas próprias do “ internetês ”, relacionadas ao uso de redes sociais e emails, como"p/ ”, “ tá ”, “ pra ”, “ pro ”, “ prum ”.
Aspectos morfossintáticos – De uma maneira geral, pode -se afirmar que há um domínio mediano das convenções morfossintáticas pelos participantes.
A seguir analisam -se alguns desvios relativos à regência e à concordância nominal e verbal.
Quanto à regência nominal e verbal, o desvio mais frequente é a falta do sinal indicativo da crase, o que revela que o usuário não tem consciência de que “ a ” reúne um artigo e uma preposição “ a ”, exigida pela regência do termo anterior ( substantivo, adjetivo ou verbo ).
Outro problema relacionado à regência verbal e à nominal é a ausência de preposição antes de pronome relativo, processo generalizado na modalidade oral da língua, em situações de registro informal.
Também é recorrente o emprego inadequado de preposição, por exemplo : “ ações a evitá -los ”.
Quanto à concordância nominal e verbal, seis aspectos merecem destaque : a ) ausência ou presença indevida de marca de plural ( com sujeito anteposto ou posposto ).
Exemplos : “ para que a distribuição de água cheguem ”, “ fica claro os benefícios ” ;"evitar que aconteça esses desastres";"criar um sistema de coletagem de ideias para que seja amenizada as perdas";"equipamentos que medi ” ;"Dessa forma corroborando para que esferas do nível federal enalteça mediadas a serem tomadas...";"para que houvessem mais áreas drenadas"; b ) aparecimento da marca de plural em verbos ou adjetivos relacionados a núcleos substantivos no singular, devido à presença de modificadores no plural.
Exemplos :"o aumento dos acidentes ambientais devem ser evitados ” ; “ a fiscalização sobre os laudos deveriam ”, evidenciando um processo de hipercorreção ; c ) ausência de acento circunflexo na forma plural do presente do indicativo do verbo “ ter ”, que é considerada como um desvio na concordância verbal e não na acentuação gráfica :"os estados e municípios tem que aplicar medidas punitivas a empresas que desmatam"; d ) ausência de concordância de gênero, no âmbito de sintagmas nominais longos, em que o adjetivo antecede ou está afastado do substantivo :"seria necessário uma reeducação";"é necessário a mobilização"; e ) emprego indevido de pronomes, reproduzindo hábito da oralidade, por exemplo, uso do pronome reto como objeto : “ deixar elas isoladas ” ; f ) confusão entre formas verbais : “ apoia -se ” por “ apoiasse".
Aspectos textuais – A maior parte das respostas foi constituída por textos estruturados, com justificativas para as propostas apresentadas, constituindo um pequeno texto de base dissertativa, mas dividido em dois parágrafos independentes, sem conexão Várias marcas de oralidade foram identificadas, embora não em alta frequência : o uso do pronome relativo “ onde ” como relativo universal, falta de artigo definido antes de substantivo, expressões informais, eliminação de preposições.
Aspectos Vocabulares – O desempenho nesta competência foi muito diversificado, revelando diferentes graus de domínio vocabular.
Como o Enade/2019 avaliou graduandos de Engenharia, pode -se observar nitidamente, em uma parte dos participantes, o domínio de um vocabulário específico na discussão da questão dos desastres ambientais e das propostas de intervenção para esse grave problema.
Por outro lado, uma grande quantidade de respostas limitou -se ao senso comum, marcado por um vocabulário pouco específico, de caráter genérico.
Quanto ao grau de formalidade do vocabulário empregado, pode -se afirmar que os participantes empregaram um registro semiformal, sem a utilização de gírias ou palavras típicas da oralidade.
Vários tipos de inadequação foram observados na seleção e na utilização do vocabulário de uma parte dos participantes : a ) excessiva repetição de certas palavras, revelando limitação de repertório vocabular ; b ) expressões da oralidade – por exemplo : “ colocar isso na cabeça das pessoas ”, “ não saem do lugar ” ; c ) seleção vocabular incompatível com o contexto, gerando falta de inteligibilidade, por exemplo : “ no requisito ao combate ” ; “ moradias dignas de baixo custo ” ; d ) falta de domínio de vocabulário mais abstrato e de maior complexidade, essencial ao desenvolvimento do texto de base dissertativa : uso de"corroborando"no lugar de"colaborando"; “ suposto ” por “ pressuposto"; “ devemos obter algumas atitudes ” ; “ recorrentes ” por “ decorrentes"; “ destacados ” por “ descartados ” ; “ eminente ” por “ iminente ” ; e ) emprego de palavras inexistentes, por exemplo : “ trabalhabilidade ”, “ dizrespeito ” ;"acessiva ”,"coletagem"; f ) imprecisão vocabular, como o uso da palavra “ coisa ” em vários textos, demonstrando o domínio precário da norma-padrão da Língua Portuguesa, com presença constante de traços de oralidade.
Aspectos textuais – A maior parte das respostas foi constituída por textos estruturados, com justificativas para as propostas apresentadas, constituindo um pequeno texto de base dissertativa, mas dividido em dois parágrafos independentes, sem conexão lógica, devido ao que foi solicitado pelo comando.
Entretanto, foram identificadas algumas estruturas textuais que não atenderam à necessidade de se construir um texto para apresentação das propostas de intervenção : a ) respostas com sintagmas nominais curtos, em geral, de duas ou três palavras.
Exemplos :"secas, inundações, alagamentos";"redução dos níveis de risco"; b ) respostas com desenvolvimento do conteúdo da palavra ou expressão nuclear ( sintagmas nominais mais amplos, em geral de mais de quatro palavras ).
Exemplos :"minimização dos impactos futuros dos desastres ambientais";"fortalecimento das resiliências ambiental, social e econômica";"contenção de encostas em áreas de risco"; c ) respostas com estrutura frasal completa, mas com pouco desenvolvimento, reduzidas a apenas um período : Exemplos : “ dragar ” os rios que cortam as grandes cidades";"estimular o trabalho voluntário".
As respostas que se constituíram como textos desenvolvidos apresentaram distintos desempenhos : textos bem elaborados, com estruturação, fluência e coerência ; textos de desempenho mediano, com alguns desvios de estruturação ; e textos bem precários, com muitos desvios e comprometimento da coerência e da fluência textual, refletindo tendências gerais na escrita de estudantes de diferentes níveis de ensino.
Seguem alguns desvios observados : a ) sequência justaposta de ideias sem encaixamentos sintáticos ; b ) redução de estruturas subordinadas, ao lado do aumento na frequência de estruturas coordenadas e absolutas ; c ) redução no uso de conectores para expressar relações lógicas essenciais à construção do texto ; d ) repetição exaustiva de termos sem a utilização de procedimentos mais sofisticados de referenciação ( hiperonímias, hiponímias, nominalizações, expressões metafóricas ) ; e ) frases fragmentadas que comprometem a estrutura lógico-gramatical, como frases formadas apenas por oração subordinada, sem oração principal.
Exemplos : - oração subordinada solta : “ No qual forneceria moradias dignas de baixo custo ” ; - frase iniciada por gerúndio solta : “ Evitando assim a ocupação de lugares impropios ( por impróprios ) para viver ” ;"Podendo também fazer eventos na intenção..."; - truncamento sintático : “... isso faz com que os bueiros fiquem entupidos e aconteça e a água...".
Quanto à utilização dos sinais de pontuação, observou -se uma grande diversidade de situações nos textos analisados, desde textos bem pontuados até textos completamente sem pontuação, inclusive sem ponto final.
São os seguintes os tipos de problemas encontrados : - emprego de vírgula para separar o sujeito e o predicado, procedimento muito frequente ; - emprego de vírgula no lugar do ponto para separar ideias que constituem períodos distintos ; - ocorrência de apenas uma das vírgulas para destacar uma palavra, uma expressão ou uma oração encaixada ; - ausência de vírgula para separar elementos de uma enumeração ; - ausência de vírgula para separar oração adjetiva explicativa ou utilização inadequada de vírgula para separar oração adjetiva restritiva ; - emprego de ponto e vírgula no lugar de vírgula ; - ausência de ponto final para \textbf{encerrar} o texto. 6.3.1.8 Análise de Língua Portuguesa da Questão Discursiva 2 do Componente de Formação Geral Os dados de Tecnologia em Estética e Cosmética, obtidos a partir das respostas à questão discursiva 2 do Componente de Formação Geral, no que tange à Língua Portuguesa, encontram -se na Tabela 6.17 e no Gráfico 6.10.
Nesse aspecto, os estudantes de todo o Brasil obtiveram Média 37,0, menor do que a obtida na questão 1.
A maior Média em relação à Língua Portuguesa foi obtida na região Sul ( 38,5 ), e a menor, na região Norte ( 33,6 ).
Quanto à variabilidade das notas, o Desvio padrão de todo o Brasil foi 27,2.
O menor Desvio padrão foi obtido na região Nordeste ( 26,9 ) e o maior Desvio padrão foi obtido na região Norte ( 28,0 ).
A Mediana das notas de Língua Portuguesa para a questão 2 foi 45,0 para o Brasil, como um todo, a mesma sendo obtida em todas as regiões.
A nota Máxima para todo o Brasil foi 90,0, com, pelo menos, um estudante tirando essa nota em todas as regiões, assim como a nota Mínima foi zero em todas as regiões do país.
No Gráfico 6.10, mostra -se a distribuição das notas de Língua Portuguesa da questão 2 do Componente de Formação Geral.
Observa -se que a maior frequência ( 25,5 \% ) corresponde à dos estudantes que que deixaram a resposta à questão 2 em branco.
Destacam -se, também, os estudantes com nota no intervalo ( 50 ; 60 ], representando 21,6 \% do total. 6.3.1.9 Comentários sobre a correção de Língua Portuguesa das respostas à Questão Discursiva 2 A Questão 2 de Formação Geral abordava o tema do desenvolvimento científico e do desenvolvimento tecnológico do Brasil, tomando como base um fragmento de texto publicado em site oficial de um órgão governamental de pesquisa na internet, o Ipea.
O texto motivador informava que o Brasil está em posição intermediária em todos os indicadores de produção e utilização de conhecimentos e novas tecnologias.
Além desse texto, a questão apresentava uma imagem que estabelecia uma ligação entre Instituição de Ensino Superior - Empresa - Setor Público.
Com base nessa reflexão inicial despertada pelo texto motivador, o comando solicitava que o participante citasse dois ganhos para o campo científico e dois ganhos para o campo econômico no país, resultantes da relação entre Instituição de Ensino Superior - Empresa - Setor Público.
Como o comando solicitava apenas uma “ citação ” de ganhos, sem referência à exigência de justificativas ou explicações que ensejariam a elaboração de um texto dissertativo sobre o tema, observou -se uma grande fragmentação do pensamento abstrato e, por consequência, uma grande fragmentação discursiva.
Essa diversidade pode ter sido consequência de uma certa insegurança entre os participantes a respeito de qual modelo de resposta produzir.
Assim, houve uma grande diversidade de modelos de resposta, desde a citação de duas palavras ou expressões, sem desenvolvimento, até a construção de respostas com detalhamento dos ganhos sugeridos, acompanhados por uma justificativa de caráter dissertativo, passando por vários modelos intermediários.
Essa diversidade tipológica, que será mais detalhada adiante, exigiu mudanças nos critérios habitualmente utilizados para avaliação de uma produção textual.
Como resultado desse enunciado, as respostas foram majoritariamente apresentadas sob a forma de dois itens, cada um atendendo à"citação"dos ganhos científicos e dos ganhos econômicos.
O que variou foi o formato textual dessa"citação", conforme detalhado a seguir. a ) Respostas com palavras ou expressões ( sintagmas nominais curtos, em geral, de uma a três palavras ), constituindo uma frase nominal com núcleo substantivo : Exemplos : “ Descobertas de algumas doenças ” ; “ Frigoríficos de frango ” ; “ Valorização ” ; “ Infraestrutura ” ; “ Laboratórios ” ; “ Campos de estudo ” ; “ Tecnologia ” ; “ Automação ”. b ) Respostas com desenvolvimento do conteúdo da palavra ou expressão nuclear ( sintagmas nominais mais amplos, em geral de mais de quatro palavras ), constituindo uma frase nominal com núcleo substantivo : c ) Respostas iniciadas por palavra ou expressão ( sintagmas nominais curtos ), mas com a presença de expressão verbal na complementação do núcleo sintático, constituindo uma frase nominal ( com núcleo substantivo ) complexa devido à presença de oração subordinada com núcleo verbal : Exemplos : “ Setor econômico possibilitando a oferta de emprego ao cidadão ” ; “ Auxílio financeiro para que os pesquisadores possam permanecer na academia ” ; “ Desenvolvimento de novas metodologias que permitam aumentar a produtividade de áreas agriculturáveis brasileiras ” ; “ Flexibilização de leis ( ou criação ) de modo a estimular a utilização de novas tecnologias ” ; “ Maior investimento estrangeiro para incentivar novas pesquisas e reconhecimento internacional ”. d ) Respostas com estrutura frasal completa ( com núcleo verbal ), mas com pouco desenvolvimento, reduzidas a apenas um período : Exemplos : “ Possibilitar a oferta de emprego ao cidadão e investimentos para as empresas, com fácil acesso para o pequeno empresário ” ; “ Investindo nesses recursos de pesquisas cientificas, obteve -se um retorno na economia do país ” e ) Respostas mais desenvolvidas, com justificativas, constituindo um pequeno texto de base dissertativa, dividido ou não em dois parágrafos independentes, sem conexão lógica : Exemplo : “ A consultoria de membros laboratoriais de instituições de ensino superior em pequenas e médias empresas pode agregar valor tanto no processo produtivo da empresa quanto numa maior qualificação desses “ \textbf{consultores} ” no mercado de trabalho. ( Maior empregabilidade ) ”.
Para avaliar essa diversidade de situações discursivas da Questão 2, não foi possível utilizar o padrão de resposta que vem sendo utilizado para a avaliação do desempenho linguístico dos participantes dos Enade anteriores e na correção da Questão 1.
Aquele padrão se revelou inadequado para a realidade das respostas analisadas na Questão 2, porque uma parcela significativa dos participantes elaborou respostas curtas, compostas por itens isolados, criando praticamente uma lista de palavras, expressões ou frases soltas, sem conexão semântica ou textual.
Exemplos : “ Aquisição de equipamentos e materiais necessários ao desenvolvimento de pesquisas ” ; “ Descobertas no ramo da medicina relacionadas à cura de novas doenças ” ; “ Descoberta de novas tecnologias pelas instituições de pesquisa ” ; “ Melhoria das instituições de ensino superior pela ajuda empresarial ” ; “ Desenvolvimento de energias sustentáveis e economicamente acessíveis ”. c ) Respostas iniciadas por palavra ou expressão ( sintagmas nominais curtos ), mas com a presença de expressão verbal na complementação do núcleo sintático, constituindo uma frase nominal ( com núcleo substantivo ) complexa devido à presença de oração subordinada com núcleo verbal : Exemplos : “ Setor econômico possibilitando a oferta de emprego ao cidadão ” ; “ Auxílio financeiro para que os pesquisadores possam permanecer na academia ” ; “ Desenvolvimento de novas metodologias que permitam aumentar a produtividade de áreas agriculturáveis brasileiras ” ; “ Flexibilização de leis ( ou criação ) de modo a estimular a utilização de novas tecnologias ” ; “ Maior investimento estrangeiro para incentivar novas pesquisas e reconhecimento internacional ”. d ) Respostas com estrutura frasal completa ( com núcleo verbal ), mas com pouco desenvolvimento, reduzidas a apenas um período : Exemplos : “ Possibilitar a oferta de emprego ao cidadão e investimentos para as empresas, com fácil acesso para o pequeno empresário ” ; “ Investindo nesses recursos de pesquisas cientificas, obteve -se um retorno na economia do país ” e ) Respostas mais desenvolvidas, com justificativas, constituindo um pequeno texto de base dissertativa, dividido ou não em dois parágrafos independentes, sem conexão lógica : Exemplo : “ A consultoria de membros laboratoriais de instituições de ensino superior em pequenas e médias empresas pode agregar valor tanto no processo produtivo da empresa quanto numa maior qualificação desses “ \textbf{consultores} ” no mercado de trabalho. ( Maior empregabilidade ) ”.
Para avaliar essa diversidade de situações discursivas da Questão 2, não foi possível utilizar o padrão de resposta que vem sendo utilizado para a avaliação do desempenho linguístico dos participantes dos Enade anteriores e na correção da Questão 1.
Aquele padrão se revelou inadequado para a realidade das respostas analisadas na Questão 2, porque uma parcela significativa dos participantes elaborou respostas curtas, compostas por itens isolados, criando praticamente uma lista de palavras, expressões ou frases soltas, sem conexão semântica ou textual.
Assim, realizou -se uma adequação dos critérios de avaliação para atender às características das respostas produzidas, em função das exigências estabelecidas pelo enunciado : “ citar ganhos ”.
A análise do desempenho linguístico esperado de um estudante que está concluindo seu curso superior e, devido a isso, está realizando o exame do Enade, tem se pautado no princípio de que ele deve utilizar seus conhecimentos sobre os dois temas de Formação Geral apresentados e estruturar seus textos, de base dissertativa, de acordo com as características do registro formal adequado à situação comunicativa – avaliação de conhecimentos.
Essa configuração determina exigências quanto : à adequação da seleção vocabular, à progressão temática coerente do conteúdo, à estruturação sintática dos períodos, à organização lógica das ideias, à utilização de procedimentos de encadeamento textual e de referenciação, à obediência às exigências morfossintáticas próprias da modalidade escrita da norma-padrão, ao respeito às regras ortográficas e às regras de acentuação gráfica.
São competências distintas, que permitem um mapeamento detalhado do domínio dos recursos disponíveis na Língua Portuguesa para a comunicação escrita formal, que constituem um acervo linguístico dos graduandos para enfrentar as exigências de sua vida profissional futura.
Na Questão 2 da edição do Enade de 2019, entretanto, a avaliação dos aspectos textuais ficou extremamente prejudicada, porque a maioria das respostas não apresentou material linguístico suficiente para que essa competência fosse pontuada.
Devido a isso, a competência relativa ao domínio dos aspectos textuais não foi avaliada na Questão 2 do Enade/2019 e, em \textbf{compensação}, foram ampliados, consequentemente, os valores relativos às demais competências : atendimento às convenções ortográficas, adequação na seleção vocabular e atendimento às convenções morfossintáticas.
Essa solução se justifica, também, porque a avaliação dos aspectos textuais criaria uma desigualdade ente os participantes : aqueles que atenderam ao comando do enunciado e produziram uma resposta composta apenas por itens para apresentar os “ ganhos ” seriam prejudicados, porque seriam penalizados por não atenderem às exigências de uma construção coerente e bem desenvolvida.
O desempenho linguístico dos participantes foi avaliado, portanto, com base nas três competências descritas a seguir : - domínio das convenções ortográficas : grafia de vogais e consoantes, uso de maiúsculas e minúsculas, emprego do hífen e acentuação gráfica ; - domínio da seleção vocabular, tendo em vista as exigências semânticas do tema desenvolvido e do registro formal da modalidade escrita culta ; - domínio das convenções de caráter morfossintático estabelecidas como modelares do ponto de vista da modalidade escrita formal culta da Língua Portuguesa : concordância nominal e verbal, regência nominal e verbal, colocação pronominal, flexão nominal e verbal, correlação entre tempos e modos verbais, ausência de marcas de oralidade.
Os critérios para análise de cada uma das competências que envolvem o domínio das convenções ortográficas, das convenções morfossintáticas e da seleção vocabular da escrita formal culta da Língua Portuguesa vêm explicitados a seguir.
Esta competência envolve :  precisão na utilização do vocabulário relacionado ao tema da questão ;  ausência de marcas de oralidade, como termos de sentido muito genérico e termos de registros mais informais.
Espera -se que o participante :  utilize repertório vocabular diversificado, sem repetição inadequada de termos ;  recorra a vocabulário preciso em relação ao tema da questão ;  complemente o sentido das palavras por meio de expressões modificadoras ( exemplo : “ Desenvolvimento de projetos ” – sentido incompleto no que diz respeito à palavra “ desenvolvimento ” e à palavra “ projetos ” ) ;  empregue um vocabulário mais formal, com ausência de marcas da oralidade, como termos de sentido muito genérico ( “ coisa ”, “ negócio ”, “ você ” ) e termos de registros mais informais ( como gírias, jargões, frases feitas, ditados populares, termos regionais ) ;  seja coerente na combinação vocabular.
O padrão de resposta aplicado para a avaliação da Questão 2 não considerou, portanto, as competências relativas aos aspectos textuais, conforme foi realizado para a Questão 1.
Ao final do processo de avaliação, ficou bem evidenciada a influência que o enunciado da Questão 2 exerceu sobre o desempenho dos participantes.
A grande diversidade de situações de resposta não expressou apenas a dificuldade de organização do pensamento sob a forma de um texto bem estruturado e coerente, mas se refletiu também nas competências relativas aos aspectos vocabulares e morfossintáticos.
Por exemplo, nas respostas mais sintéticas, constituídas por itens preenchidos por sintagmas nominais simples ( compostos apenas por um núcleo substantivo ), os estudantes apresentaram fraco desempenho nos aspectos vocabulares ao empregarem termos muito vagos ou imprecisos semanticamente ( exemplo : “ Hierarquia ” ).
Quando o núcleo desse modelo de resposta era preenchido por substantivos deverbais ( relacionados semanticamente a um verbo ), ficavam incompletos do ponto de vista morfossintático pela ausência de um complemento nominal ( exemplo : “ Valorização ” de quê? ).
Com base na avaliação das respostas, pode -se considerar que o desempenho dos participantes foi de nível mediano, em sua maior parte, com vários desvios nas três competências analisadas.
As respostas que foram avaliadas como fracas foram de dois tipos principais : respostas itemizadas, com sintagmas reduzidos a um núcleo, acompanhado ou não por modificadores ; ou respostas itemizadas constituídas por pequenos parágrafos, com muitos desvios nas três competências, principalmente na seleção vocabular e no domínio das convenções morfossintáticas.
Em comum nos dois tipos, a pouca ou nenhuma articulação entre os itens ‘ a ’ e ‘ b ’, propostos na questão.
As respostas que foram avaliadas como medianas também foram itemizadas, mas cada item era abordado por parágrafos com uma estruturação textual relativamente adequada e desvios nas três competências.
As respostas que foram avaliadas como boas apresentavam desenvolvimento textual fluente e coerente, vocabulário formal e técnico, domínio das convenções morfossintáticas e ortográficas, mesmo que apresentassem eventuais desvios em um ou mais dos aspectos avaliados.
Apresentam -se, a seguir, algumas observações sobre o desempenho dos participantes em relação às três competências avaliadas.
Aspectos ortográficos – Assim como foi observado na Questão 1, o desempenho dos participantes revelou uma diferença muito grande nos dois aspectos analisados : baixo índice de desvios de grafia padrão e grande índice de desvios de acentuação.
Em vários casos, ocorre ausência quase completa de acentuação gráfica.
Com base na análise das respostas, conclui -se, portanto, que, de uma maneira geral, os participantes dominam as convenções relativas à grafia das palavras, mas desconhecem, ou desprezam, as convenções relativas à acentuação gráfica.
Seguem exemplos de algumas ocorrências : a ) eliminação do acento indicador da sílaba tônica : - palavras proparoxítonas : “ praticos ” ( por “ práticos ” ), “ tecnologicos ” ( por “ tecnológicos ” ) ; - palavras paroxítonas : “ diminuissem ” ( por “ diminuíssem ” ), “ relevancia ” ( por “ relevância ” ) ; “ inimaginaveis ” ( por “ inimagináveis ” ) ; b ) uso indevido do acento gráfico, \textbf{localizado} em sílaba átona : “ cancêr ” ( por “ câncer ) ; c ) omissão da cedilha, do til ou do traço da letra"t": “ populacão ” ( por “ população ” ), “ especulaçoes ” ( por “ especulações ” ), “ orientacoes ” ( por “ orientações ” ) ; d ) desvios de grafia :"concientizar"( por"conscientizar"), “ tramisiveis ” ( por “ transmissíveis ” ), “ pulblico ” ( por “ público ” ) ; e ) uso indevido de inicial maiúscula : “ Periódicos ” ; f ) uso indevido ou ausência de hífen : “ micro-empresário ” ( por “ microempresário ” ), “ recém formados ” ( por “ recém-formados ” ) ; g ) ao contrário do que se esperava, foram poucas as ocorrências de abreviaturas próprias do “ internetês ”, relacionadas ao uso de redes sociais e emails, como"p/ ”, “ tá ”, “ pra ”, “ pro ”, “ prum ”.
Aspectos Vocabulares – O desempenho nesta competência foi muito diversificado, revelando diferentes graus de domínio vocabular.
A seleção e a adequação vocabular foram consideradas essenciais na avaliação das respostas itemizadas, porque a tendência dos participantes foi a de utilizar palavras muito vagas e genéricas.
Quanto ao grau de formalidade do vocabulário empregado, pode -se afirmar que os participantes utilizaram um registro semiformal, sem a utilização de gírias ou palavras típicas da oralidade.
Vários tipos de inadequação foram observados na seleção e na utilização do vocabulário de uma parte dos participantes : a ) desvios na precisão vocabular : “ Começar a produzir ”, sem a delimitação do que ‘ produzir ” ; “ o custo seria inferior ”, sem a delimitação do que tem custo ; “ Com isso o custo seria inferior ao comprador, onde iria igualar aos países mais Desenvolvidos ”, sem delimitar o que se igualaria aos países mais desenvolvidos ; b ) seleção vocabular inadequada : “ Estudantes que estão cada vez mais inteligentes e com grandes ideias afim de \textbf{ajudar} essa empresa ” ; “ recursos de supra importância ” ;"agregando valor ao nome brasileiro ao exterior ” ; “ os profissionais qualificados através de instituições superiores, promovem desempenho econômico abrangente à economia do país, pois a geração de empregos e presença de concursos públicos privados, geram saldo positivo em contratações e giro da economia promovida pelas empresa do setor público ” ; c ) neologismos não \textbf{autorizados} pelos dicionários : “ tecnologias sustentais ” ; “ Ingressão no meio empresarial ”.
Aspectos morfossintáticos – De uma maneira geral, pode -se afirmar que há um domínio mediano das convenções morfossintáticas pelos participantes, embora a estrutura itemizada das respostas, ocorrida na Questão 2, tenha favorecido muitos desvios morfossintáticos : a ) falta de paralelismo entre elementos que introduzem respostas itemizadas : “ maior número de doutores formados pelas universidades ” ; “ maior reconhecimento mundial relacionado a educação ” ; “ maior investimento de multinacionais, gerando empregos ” ; “ estimular o empreendimento interno ” ( no lugar de “ estímulo ” para manter o paralelismo com os outros substantivos dos itens anteriores ). b ) falta de preposição para manter o paralelismo em enumeração : “ maior geração de conhecimento e ( de ) informação ” ; “ maior disponibilização de recursos e ( de ) dados ” ; “ relatórios de realização e ( de ) atividade econômica ” ; “ a produção qualificada da ciência e ( da ) tecnologia ” ; “ \textbf{ajudar} na qualidade e ( na ) rapidez da produção ”. c ) falta de artigo na enumeração de itens : “ a educação e ( o ) conhecimento elevam uma sociedade ” ; “ onde é melhor realizar melhorias para as cidades, ( os ) estados e ( o ) país ” ; “ isso volta como serviço para as empresas e ( o ) setor público ” ; “ juntamente com o setor privado e ( o ) público ” ; “ criação de novas tecnologias para as empresas e ( o ) setor público ” ; “ desenvolver parceria entre as universidades e ( as ) empresas privadas ” ; “ investimentos viabilizam ( a ) produção e ( a ) venda de tecnologia ” ; “ profissionais capacitados para os trabalhos e ( as ) pesquisas ” ; “ com a contribuição e ( a ) colaboração dos setores públicos e empresariais ” ; “ surgimento das novas ideias, ( das ) pesquisas e ( dos ) projetos ” ; “ o preço cobrado pelos produtos e ( pelas ) empresas seria menor ” ; “ novas descobertas nos campos da saúde e ( da ) engenharia ”. d ) ausência de complementos nominais junto a substantivos abstratos deverbais, principalmente nas respostas itemizadas : “ valorização ”, no lugar de “ valorização de... ” ; “ patrocínio ” no lugar de “ patrocínio de... ”.
Desvios de regência nominal e verbal :  o desvio mais frequente é a falta do sinal indicativo da crase, o que revela que o usuário não tem consciência de que o termo “ a ” reúne um artigo"a"e uma preposição “ a ”, exigida pela regência do termo anterior ( substantivo, adjetivo ou verbo ) : “ maior reconhecimento mundial relacionado a ( por à ) educação ” ;  ausência de preposição antes de pronome relativo, processo generalizado na modalidade oral da língua, em situações de registro informal : “ as instituições de ensino superior ("em") que as empresas investem ” ;  emprego inadequado de preposição : “ o custo seria inferior ao ( por"para o") comprador ”.
Desvios de concordância nominal e verbal :  ausência ou presença indevida de marca de plural ( com sujeito anteposto ou posposto ).
Exemplos : “ Os setores público e ensino superior poderiam desenvolverem ( por “ desenvolver ” ) pesquisas para empresas seja ( por “ sejam ” ) elas do tipo tecnologico ( por “ tecnológico ” ), saúde, transporte entre outros ” ; “ já as empresas poderiam financiar projetos de pesquisas para os setores público e ensino superior que também poderia se utilizarem ( por “ utilizar ” ) de suas partes estruturais e acadêmicas para desenvolvimento tecnologico ” ( por “ tecnológico ” ) ; “ as empresas que produz ( por “ produzem ” ) os materiais ” ; “ o engajamento entre ambos elevam ( por “ eleva ” ) a economia ”.  Aparecimento da marca de plural em verbos ou adjetivos relacionados a núcleos substantivos no singular devido à presença de modificadores no plural.
Exemplos :"o financiamento das bolsas de estudo devem ( por “ deve ” ) ser ampliados ” ( por"ampliado") ; “ o acesso aos cursos deveriam ( por “ deveria ” ) ser incentivados ( por"incentivado") pelas empresas ”, evidenciando um processo de hipercorreção.  Ausência de acento circunflexo na forma plural do presente do indicativo dos verbos “ ter ” e “ vir ”, considerada como um desvio na concordância verbal e não na acentuação gráfica.
Exemplos : “ ganhos na economia vem ( por “ vêm ” ) através dos aumentos com exportação ” ;"as universidades públicas tem ( por"têm") recebido menor quantidade de verbas para suas atividades".  Ausência de concordância de gênero, no âmbito de sintagmas nominais longos, em que o adjetivo antecede ou está afastado do substantivo.
Exemplos :"seria necessário uma reeducação"( por"necessária") ;"é necessário ( por"necessária") a mobilização";  Emprego de “ onde ” não locativo.
Exemplo :  emprego inadequado de preposição : “ o custo seria inferior ao ( por"para o") comprador ”.
Desvios de concordância nominal e verbal :  ausência ou presença indevida de marca de plural ( com sujeito anteposto ou posposto ).
Exemplos : “ Os setores público e ensino superior poderiam desenvolverem ( por “ desenvolver ” ) pesquisas para empresas seja ( por “ sejam ” ) elas do tipo tecnologico ( por “ tecnológico ” ), saúde, transporte entre outros ” ; “ já as empresas poderiam financiar projetos de pesquisas para os setores público e ensino superior que também poderia se utilizarem ( por “ utilizar ” ) de suas partes estruturais e acadêmicas para desenvolvimento tecnologico ” ( por “ tecnológico ” ) ; “ as empresas que produz ( por “ produzem ” ) os materiais ” ; “ o engajamento entre ambos elevam ( por “ eleva ” ) a economia ”.  Aparecimento da marca de plural em verbos ou adjetivos relacionados a núcleos substantivos no singular devido à presença de modificadores no plural.
Exemplos :"o financiamento das bolsas de estudo devem ( por “ deve ” ) ser ampliados ” ( por"ampliado ” ) ; “ o acesso aos cursos deveriam ( por “ deveria ” ) ser incentivados ( por"incentivado") pelas empresas ”, evidenciando um processo de hipercorreção.  Ausência de acento circunflexo na forma plural do presente do indicativo dos verbos “ ter ” e “ vir ”, considerada como um desvio na concordância verbal e não na acentuação gráfica.
Exemplos : “ ganhos na economia vem ( por “ vêm ” ) através dos aumentos com exportação ” ;"as universidades públicas tem ( por"têm") recebido menor quantidade de verbas para suas atividades".  Ausência de concordância de gênero, no âmbito de sintagmas nominais longos, em que o adjetivo antecede ou está afastado do substantivo.
Exemplos :"seria necessário uma reeducação"( por"necessária") ;"é necessário ( por"necessária") a mobilização";  Emprego de “ onde ” não locativo.
Exemplo : “ Com isso o custo seria inferior ao comprador, onde iria igualar aos países mais desenvolvidos ”. 6.3.1.10 Considerações finais sobre a correção de Língua Portuguesa das respostas às Questões Discursivas do Componente de Formação Geral A análise do desempenho linguístico das duas questões de Formação Geral do Enade/2019 permite chegar a algumas conclusões que podem colaborar para um aperfeiçoamento dos cursos superiores e para uma adequação do exame aos seus propósitos institucionais.
Um desses propósitos é a avaliação sobre a competência discursiva dos graduandos, conforme determinação do Inep : as questões discursivas devem fornecer material para que as respostas sejam avaliadas em aspectos como clareza, coerência, coesão, estratégias argumentativas, utilização de vocabulário adequado e correção gramatical do texto.
Com base na constatação de que nossos estudantes têm um domínio apenas mediano dessas exigências de construção de um texto formal culto, conclui -se que é preciso aperfeiçoar sua formação superior, dando -lhes acesso às ferramentas necessárias para que sejam bem-sucedidos nas diferentes situações de sua vida cidadã – profissional, cultural e política.
Essa é uma contribuição efetiva que os resultados do Enade podem oferecer às instituições de ensino superior do país.
Um segundo aspecto que ficou evidenciado a partir da análise das respostas foi o pouco conhecimento dos graduandos a respeito de temas exteriores ao âmbito específico de sua profissão, ligados à realidade nacional e global e a outras áreas do conhecimento, objetivo das questões de Formação Geral do Enade.
No exame de 2019, isso ficou evidenciado na diferença de desempenho entre estudantes de áreas de Engenharia e de outras áreas de formação na Questão 1, que tratava de desastres ambientais.
Mesmo sem ter acesso à sua área de formação, ao se analisarem as respostas, evidenciou -se o domínio de um vocabulário especializado por uma parte dos graduandos, enquanto a outra parte limitou -se ao emprego de um vocabulário próprio do senso comum, com palavras genéricas e vagas.
Por exemplo, se o tema tivesse sido “ pandemia ”, os estudantes das áreas da saúde teriam, provavelmente, tido um desempenho vocabular distinto, utilizando termos técnicos desconhecidos pelos estudantes de Engenharia.
Outro ponto a ser observado, já em consequência da análise da Questão 2, é o baixo nível de conhecimento e de reflexão dos estudantes sobre um aspecto tão importante da vida nacional : a relação entre o público e o privado, no que diz respeito ao ensino superior, à pesquisa universitária e seu papel no desenvolvimento brasileiro. 6.3.2 Componente de Conhecimento Específico Na parte da prova relativa às questões discursivas do Componente de Conhecimento Específico ( Tabela 6.18 ), observa -se que a Média foi mais alta do que para as questões discursivas do Componente de Formação Geral.
Enquanto no Componente de Formação Geral a Média para estudantes de Tecnologia em Estética e Cosmética de todo o Brasil foi 21,3, na parte de Conhecimento Específico, a Média foi 31,6.
A maior Média deste componente foi obtida pelos estudantes da região Sudeste ( 32,9 ), e a menor, pelos da região Norte ( 27,2 ).
Nas demais regiões as Médias foram ; 31,4, na região Nordeste ; 31,6, na região Sudeste ; e 30,6, na Centro-Oeste.
Quanto à variabilidade das notas, o Desvio padrão de todo o Brasil foi 14,8.
O maior Desvio padrão foi encontrado na região Norte ( 16,2 ), e o menor, na região Centro-Oeste ( 13,0 ).
A maior nota Máxima, 91,7, foi obtida na região Sudeste.
Nas demais regiões, as notas máximas foram : 73,3, na região Norte, 78,3, na região Nordeste ; 85,0, na região Sul ; e 66,7, na região Centro-Oeste.
A nota Mínima ( 0,0 ) foi obtida por, pelo menos, um estudante em todas as regiões do Brasil.
A Mediana do Brasil, como um todo, foi 31,7, a mesma das regiões Sudeste e Centro-Oeste.
Nas regiões Nordeste e Sul a Mediana foi maior ( 33,3 ) e na região Norte foi menor ( 25,0 ).
A respeito da adequação do exame aos seus propósitos institucionais, vale observar que a Questão 2 não atendeu às exigências de avaliação, tanto do domínio discursivo dos estudantes, quanto do seu conhecimento sobre temas diversos ligados à realidade nacional e à global, assim como a outras áreas do conhecimento exteriores ao âmbito específico de sua profissão.
Conforme foi detalhadamente explicado no corpo deste relatório, o enunciado da Questão 2 induziu os estudantes a criarem respostas vagas e insatisfatórias, com repertório lexical pouco diversificado, impreciso ou sem desenvolvimento, com tessitura textual comprometida por truncamentos sintático-semânticos, reflexo da fragmentação do pensamento.
Se o objetivo é avaliar a competência discursiva dos formandos brasileiros, a elaboração de um texto dissertativo certamente ofereceria material relevante para essa análise, por incentivar a reflexão e a organização do pensamento em uma sequência textual coerente e complexa. 6.3.2 Componente de Conhecimento Específico Na parte da prova relativa às questões discursivas do Componente de Conhecimento Específico ( Tabela 6.18 ), observa -se que a Média foi mais alta do que para as questões discursivas do Componente de Formação Geral.
Enquanto no Componente de Formação Geral a Média para estudantes de Tecnologia em Estética e Cosmética de todo o Brasil foi 21,3, na parte de Conhecimento Específico, a Média foi 31,6.
A maior Média deste componente foi obtida pelos estudantes da região Sudeste ( 32,9 ), e a menor, pelos da região Norte ( 27,2 ).
Nas demais regiões as Médias foram ; 31,4, na região Nordeste ; 31,6, na região Sudeste ; e 30,6, na Centro-Oeste.
Quanto à variabilidade das notas, o Desvio padrão de todo o Brasil foi 14,8.
O maior Desvio padrão foi encontrado na região Norte ( 16,2 ), e o menor, na região Centro-Oeste ( 13,0 ).
A maior nota Máxima, 91,7, foi obtida na região Sudeste.
Nas demais regiões, as notas máximas foram : 73,3, na região Norte, 78,3, na região Nordeste ; 85,0, na região Sul ; e 66,7, na região Centro-Oeste.
A nota Mínima ( 0,0 ) foi obtida por, pelo menos, um estudante em todas as regiões do Brasil.
A Mediana do Brasil, como um todo, foi 31,7, a mesma das regiões Sudeste e Centro-Oeste.
Nas regiões Nordeste e Sul a Mediana foi maior ( 33,3 ) e na região Norte foi menor ( 25,0 ).
No Gráfico 6.11, representa -se a distribuição das notas nas questões discursivas no Componente de Conhecimento Específico.
A moda dessa distribuição ocorre no intervalo ( 30 ; 40 ], com 26,5 \% de frequência.
Observa -se que no intervalo [ 0 ; 10 ], com 8,9 \% do total de participantes, estão computadas as ocorrências de respostas em branco para as três questões discursivas de Conhecimentos Específicos.
A análise de cada uma destas questões será feita a seguir. 6.3.2.1 Análise da Questão Discursiva 3 do Componente de Conhecimento Específico Na questão 3, cujos resultados aferidos se encontram descritos na Tabela 6.19, a Média dos estudantes de todo o Brasil foi 11,8.
O desempenho dos estudantes nessa questão foi inferior ao verificado nas questões 4 e 5.
A menor Média nessa questão foi obtida pelos estudantes da região Centro-Oeste ( 6,4 ), enquanto a maior Média foi obtida na região Norte ( 16,5 ).
Quanto à variabilidade das notas, o Desvio padrão de todo o Brasil foi 14,5.
O maior Desvio padrão foi obtido na região Norte ( 16,4 ), enquanto o menor foi obtido na região Centro-Oeste ( 10,7 ).
A nota Máxima, 100,0 pontos, foi alcançada por, pelo menos, um estudante das regiões Sudeste e Sul.
A menor nota Máxima foi obtida nas regiões Nordeste e Centro-Oeste ( 75,0 ) e na região Norte a nota Máxima foi 90,0.
A Mediana do Brasil, como um todo, foi 10,0, a mesma para as regiões Nordeste, Sudeste e Sul.
Na região Norte, a Mediana foi 15,0, e, na região Centro-Oeste, foi zero ( 0,0 ).
A nota Mínima ( 0,0 ) foi a mesma em todas as regiões do Brasil.
No Gráfico 6.12, mostra -se a distribuição das notas na questão discursiva 3, do Componente de Conhecimento Específico da área de Tecnologia em Estética e Cosmética.
A moda dessa distribuição, com 38,0 \% do total de participantes, é a categoria zero, dos estudantes que, tendo respondido à questão 3, receberam nota zero.
Observa -se uma concentração de notas até 20, que chega a 83,1 \% do total de participantes. 6.3.2.2 Comentários sobre as respostas à Questão Discursiva 3 O enunciado da Questão 3 estava claro e apresentava uma boa linguagem e adequação às Diretrizes Curriculares Nacionais do curso de Tecnologia em Estética e Cosmética.
Contudo, o tema era complexo e o grau de dificuldade foi considerado como médio.
Foram observadas interpretações diferentes ao que o enunciado propunha, a parte “ explique de que forma a drenagem linfática ” para muitos foi interpretada como execução, diferentemente do que intencionava o enunciado, que o traduzia como ação fisiológica.
Dessa forma, foi a questão com o menor desempenho.
Observou -se que os estudantes, em suas respostas, não explicaram a ação fisiológica e o objetivo do uso da Drenagem Linfática Manual ( DLM ) e sua ação sobre edema.
Outra abordagem frequente foi sobre a aplicação da técnica da DLM, seus movimentos de execução e da técnica de drenagem linfática reversa aplicada na prática clínica e que correlacionava com o caso clínico da questão, mas não respondia ao que foi solicitado.
Durante a correção da questão, observou -se que a maioria dos estudantes respondeu somente sobre a relação ao edema e equimoses, por conseguinte não respondendo sobre o processo de cicatrização e hipoestesia.
Muitos estudantes abordaram em suas respostas a relação da execução da técnica, com menção à técnica reversa e à forma de aplicação das manobras, não abordando a parte fisiológica – que era o perguntado – e a ação da drenagem linfática manual.
O desempenho dos estudantes, de forma geral, foi muito fraco.
Mesmo sendo uma questão que retratava um tema importante para o tecnólogo em estética e cosmética e que apresentava um enunciado claro e objetivo, percebeu -se uma grande dificuldade no entendimento do mesmo.
No geral, o desempenho dos estudantes na Questão 3 pode ser considerado muito fraco.
De 5.286 participantes, 173 deixaram a resposta a essa questão em branco, outros 59 tiveram suas respostas consideradas como nulas, desconsideradas ou, ainda, eram respostas sobre tema de outra questão da prova ( resposta trocada ).
As demais respostas foram classificadas em três grupos : fracas, medianas e boas, como discriminado a seguir.
A maioria das respostas, um total de 4.470, foi classificada como fraca por receber notas de zero a 30.
A quantidade de respostas que com notas de 31 a 70 foi 555, um grupo bem mais reduzido.
Já as 29 respostas que receberam notas maiores do que 71 foram consideradas boas, e dessas, quatro obtiveram a nota máxima 100.
A questão foi a que obteve o pior rendimento entre as questões discursivas.
O tema da questão exigiu um embasamento fisiológico da técnica e sua ação no pós-operatório.
A correção dessa questão discursiva revelou uma possível lacuna na formação do tecnólogo em estética e cosmética, na fundamentação da técnica, suas bases fisiológicas e ações terapêuticas.
Foi observada também uma dificuldade de interpretação por parte dos estudantes. 6.3.2.3 Análise da Questão Discursiva 4 do Componente de Conhecimento Específico Na Tabela 6.20, constam as informações relativas à questão 4 do conjunto de questões do Componente de Conhecimento Específico.
O desempenho dos estudantes nessa questão ( Média 40,2 ) foi bem superior ao verificado na questão 3 ( Média 11,8 ) e um pouco inferior ao verificado na questão 5 ( Média 42,8 ).
A maior Média foi registrada na região Nordeste ( 45,6 ), e a menor, na região Norte ( 31,9 ).
No Gráfico 6.13, está representada a distribuição de notas da questão discursiva 4, do Componente de Conhecimento Específico.
Essa distribuição tem moda no intervalo ( 50 ; 60 ], com frequência 20,1 \%.
Na questão 4 foi baixo o percentual de respostas em branco ( 3,5 \% ) e de notas zero para respostas dadas pelos estudantes ( 2,3 \% ) A nota Máxima ( 100,0 ) foi atingida por, pelo menos, um concluinte de cada região.
Também foi zero a nota Mínima de todas as regiões.
A Mediana em todo o Brasil foi 40,0, o mesmo valor obtido nas regiões Sudeste e Sul.
Na região Nordeste a Mediana foi maior ( 45,0 ) e foi menor nas regiões Norte ( 30,0 ) e Centro-Oeste ( 35,0 ).
No Gráfico 6.13, está representada a distribuição de notas da questão discursiva 4, do Componente de Conhecimento Específico.
Essa distribuição tem moda no intervalo ( 50 ; 60 ], com frequência 20,1 \%.
Na questão 4 foi baixo o percentual de respostas em branco ( 3,5 \% ) e de notas zero para respostas dadas pelos estudantes ( 2,3 \% ) 6.3.2.4 Comentários sobre as respostas à Questão Discursiva 4 O enunciado da Questão 4 apresentava clareza, uma boa linguagem e um bom texto inicial que subsidiou a resposta.
O tema da questão faz parte da área de terapia capilar, abordando a disfunção de dermatite seborreica.
Solicitou -se que se abordasse a aplicação do óleo essencial de melaleuca e da argila, e também sobre o mecanismo de ação do equipamento de alta frequência como propostas de tratamento.
Esses tratamentos são considerados básicos na formação do futuro tecnólogo em estética e cosmética ; por isso, o grau de dificuldade da questão discursiva foi considerado baixo.
O padrão de respostas foi apresentado de forma detalhada.
Como o nível de dificuldade da questão era baixo, as respostas dos estudantes em geral não divergiram do padrão de resposta.
O comando da questão estava dividido em dois itens.
Durante a correção, no item ‘ a ’, de acordo com as respostas dos estudantes, não foi observada nenhuma dificuldade para responder ao que se pedia no comando.
Contudo, no item ‘ b ’, observou -se que os respondentes abordaram a ação clínica da alta frequência, mas não conseguiram explicar sua ação fisiológica.
Vale ressaltar que o aparelho de alta frequência é básico na formação do tecnólogo em estética e cosmética, sendo muito usado nas terapias faciais, com a mesma ação fisiológica.
O comando da Questão 4 estava claro e objetivo, com grau de dificuldade baixo, com tema voltado para área de terapia capilar.
O desempenho dos estudantes foi mediano, com nota média 4,2, considerando -se as respostas válidas.
No geral, o desempenho dos estudantes na Questão 4 pode ser considerado mediano.
De 5.286 participantes, 185 deixaram a resposta a essa questão em branco, outros 39 tiveram suas respostas consideradas como nulas, desconsideradas ou, ainda, eram respostas sobre tema de outra questão da prova ( resposta trocada ).
As demais respostas foram classificadas em três grupos : fracas, medianas e boas, como discriminado a seguir.
O maior grupo, 3.064 respostas, foram classificadas como medianas por receberem notas de 31 a 70.
A quantidade de respostas fracas, com notas de zero a 30, foi de 1.784, um grupo menos expressivo.
Já as 214 respostas que receberam notas maiores do que 71 foram consideradas boas, e dessas, 92 obtiveram a nota máxima 100.
O grau de dificuldade dos estudantes para responder à questão foi alto, o que talvez demonstre a necessidade de intensificar/aprofundar os conhecimentos na área de terapia capilar, objeto do tema da questão. 6.3.2.5 Análise da Questão Discursiva 5 do Componente de Conhecimento Específico Na Tabela 6.21, constam as informações relativas à questão 5 do conjunto do Componente de Conhecimento Específico.
O desempenho dos estudantes de todo o Brasil nesta questão foi o mais alto entre as três questões discursivas desse componente.
A nota Média dos estudantes de todo o Brasil foi 42,8.
A maior Média foi registrada na região Centro-Oeste ( 49,6 ), enquanto a menor Média foi registrada na região Norte ( 33,3 ).
Quanto à variabilidade das notas, o Desvio padrão dos estudantes do Brasil, como um todo, foi 25,1.
Enquanto o maior desvio foi encontrado nas regiões Norte e Nordeste ( 26,0 ), o menor foi encontrado na região Sul ( 23,2 ).
A Mediana para o Brasil e em todas as regiões foi 45,0.
Para o conjunto de estudantes de Tecnologia em Estética e Cosmética do Brasil, a nota Máxima foi 100,0, obtida por, pelo menos, um estudante de quatro regiões.
A nota Máxima foi menor na região Centro-Oeste ( 90,0 ).
A nota Mínima do Brasil foi zero, o mesmo valor encontrado em todas as regiões.
No Gráfico 6.14, é representada a distribuição das notas da questão discursiva 5 do Componente de Conhecimento Específico.
Essa distribuição tem moda no intervalo ( 50 ; 60 ], correspondendo a cerca de 18,8 \% dos respondentes.
Os participantes que tiraram zero ou que deixaram a questão em branco representam 11,2 \% do total. 6.3.2.6 Comentários sobre as respostas à Questão Discursiva 5 A Questão 5 está inserida na grande área de estética facial no subtema envelhecimento.
Esse tema é muito oportuno e básico nos cursos de Tecnologia em Estética e Cosmética.
Dessa forma, o grau de dificuldade da questão era baixo.
O comando estava claro, com linguagem adequada, além de apresentar um bom infográfico para subsidiar as respostas dos estudantes.
O padrão de resposta era adequado e abrangente.
Os estudantes, quando acertaram, pouco se afastavam do padrão.
No item ‘ a ’ observou -se que muitos estudantes citaram os dois tipos de envelhecimento, mas não os correlacionaram com os fatores desencadeantes.
Foram encontradas também as seguintes respostas : Intrínseco ( envelhecimento de dentro para fora ) e extrínseco ( de fora para dentro ), sem, contudo, detalhar os fatores desencadeantes.
No item ‘ a ’ não houve divergência entre a maioria das respostas dos estudantes e o padrão de resposta.
Os estudantes redigiram respostas similares ao padrão de resposta, na sua maioria.
Em relação ao item ‘ b ’, os estudantes elencaram modificações estruturais cutâneas decorrentes do processo de envelhecimento diferentes daquelas listadas no padrão de resposta.
Ou seja, além do que constava no padrão de resposta, foram encontradas outras modificações que não constavam no padrão.
O comando da questão estava claro e objetivo, com grau de dificuldade baixo, com tema voltado para área de Estética Facial.
No geral, o desempenho dos estudantes na Questão 5 também pode ser considerado mediano.
De 5.286 participantes, 309 deixaram a resposta a essa questão em branco, outros 22 tiveram suas respostas consideradas como nulas, desconsideradas ou, ainda, eram respostas sobre tema de outra questão da prova ( resposta trocada ).
As demais respostas foram classificadas em três grupos : fracas, medianas e boas, como discriminado a seguir.
O maior grupo, 2.757 respostas, foram classificadas como medianas por receberem notas de 31 a 70.
A quantidade de respostas fracas, com notas de zero a 30, foi de 1.548, um grupo menos expressivo.
Já as 650 respostas que receberam notas maiores do que 71 foram consideradas boas, e dessas, 37 obtiveram a nota máxima 100.
A Questão 5 foi a de melhor desempenho entre os participantes do Enade/2019, evidenciando, talvez, que o tema tenha sido bem assimilado durante a formação acadêmica dos estudantes.
Entretanto, notou -se também que alguns estudantes tiveram um pouco de dificuldade na interpretação do enunciando. 6.3.2.7 Considerações Finais sobre as questões discursivas do Componente de Conhecimentos Específicos O conteúdo das três questões discursivas estava condizente com a formação mínima necessária do tecnólogo em estética e cosmética, questões essas bem formuladas, claras, objetivas e com temas importantes para a formação desse profissional.
Os estudantes pareceram se empenhar em resolver as questões, pois na questão com o maior volume de respostas deixadas em branco esse percentual foi inferior a 6 \%.
Por outro lado, esse empenho não se refletiu em notas elevadas.
A Questão 3 foi a que obteve o pior rendimento entre as questões discursivas da parte de Conhecimentos Específicos da prova.
O tema solicitou do estudante um embasamento fisiológico da técnica de Drenagem Linfática Manual ( DLM ) e sua ação no pós-operatório.
A correção dessa questão revelou uma possível lacuna na formação do Tecnólogo em Estética e Cosmética, na fundamentação da técnica, suas bases fisiológicas e ações terapêuticas.
Pois, alguns estudantes interpretaram de forma equivocada o enunciado, respondendo como se executa a DLM, em vez de discorrer acerca da ação fisiológica da DLM.
Notou -se também que, talvez, se faça necessário um reforço na fundamentação técnica, nas bases fisiológicas e nas ações terapêuticas da DLM.
A Questão 4 foi a que obteve rendimento melhor que a discursiva 3 e pior que a discursiva 5.
Apesar de o grau de dificuldade da questão 4 ser baixo, a nota média foi razoável.
Em relação à essa questão, aqui também, talvez, se faça necessária uma maior intensificação e/ou aprofundamento nos conhecimentos da área de terapia capilar.
A Questão 5 foi a que obteve o melhor rendimento entre as questões discursivas, a de melhor desempenho entre os acadêmicos.
Evidenciando, talvez, que o tema tenha sido bem assimilado durante a formação acadêmica dos estudantes.
Entretanto, notou -se também que alguns estudantes tiveram um pouco de dificuldade na interpretação do enunciando.

% matched lemmas: ajudar, assessora, autorizar, beneficiário, brusco, catástrofe, causa-efeito, compensação, compreensível, concretizar, consultor, encerrar, formativo, julgar, laranja, localizar
\end{textsample}
